% Modernised reading — fonds Alexandre Grothendieck, shelfmark 23, pages 1-98.
%
% This is an INTERPRETATION, not a transcription. It re-reads the pages in
% today's notation and vocabulary, states the hypotheses the manuscript leaves
% implicit, and drops the critical apparatus so that the mathematics reads
% straight through. Everything the brackets and underlines carried moves into a
% footnote. It is held to being mathematically correct as it stands; where the
% manuscript is loose or mistaken, the true statement is what appears here and
% the footnote says what the page has. In any dispute of substance the
% transcription governs: cite batch-01..05.fr.tex.
%
% Pass: Opus 5 (claude-opus-5), 2026-09-21 — first pass, unchecked against the
% pages by a human. The model was chosen explicitly for this pass.
%
% Provenance of the source transcriptions. The five batches of this folder were
% not all read by the same model: batches 1, 3, 4 and 5 were transcribed by
% Fable 5.1 (claude-fable-5-1) and batch 2 by Opus 5 (claude-opus-5), all on
% 2026-09-21. Each batch header records its own. Nothing here depends on the
% difference, and nothing in this repository says how the two models differ on
% this task; the fact is recorded because a reader of the folder's apparatus is
% entitled to know that its critical judgements come from two hands.
%
% Scope: one document for the shelfmark, its five batches read as one folder.
% Eighty-one of the ninety-eight leaves carry something. Eleven are cover
% leaves bearing a title and nothing else (1, 2, 7, 17, 33, 38, 49, 54, 65, 81,
% 90); six are blank (6, 41, 42, 47, 68, 71, 85, 94 — eight, in fact); one
% (37) is a typescript foreign to the folder and is not read.
%
% SPINE. The folder is the working file of the compléments to EGA IV, and its
% cover — page 1, in his hand — says which compléments: « EGA IV, Calcul diff.
% (compléments) ». It runs in three movements, and the folder's own eleven
% cover leaves are its table of contents.
%
%   A. pp. 3-40   HE BUILDS THE INFINITESIMAL CALCULUS OF §16-17, and builds
%                 it four times over, in four registers, always reaching the
%                 SAME SHAPE: a set of objects is a principal homogeneous
%                 space under a sheaf of infinitesimal automorphisms, and that
%                 sheaf is a module of derivations.
%                   pp. 8-16   ALGEBRA. exp and log carry the primitive
%                              elements of a bigebra onto the group-like ones;
%                              a derivation exponentiates to an automorphism;
%                              the Λ-automorphisms of A ⊗_k Λ inducing the
%                              identity ARE the k-derivations A → A ⊗_k I(Λ),
%                              with Campbell-Hausdorff for group law.
%                   pp. 18-22  SCHEMES. Is^{X_0}_T(ξ,X) ≅ S(X̄/X_0): the
%                              deformations of X_0 over an infinitesimal base,
%                              trivialised on X_0, are a torsor under
%                              Aut^{X_0}_T(X).
%                   pp. 23-29  JETS. P^{(n)}(Y/X/S) and the connexions of
%                              order n, again a torsor, under the order-n
%                              infinitesimal automorphisms of the fibres.
%                   pp. 30-32  MORPHISMS. The same for a couple (u,v) over f.
%                 Around them: the adjunction formula of pp. 3-5, the two
%                 characteristic-p counterexamples of pp. 34-36 that bound
%                 §17, and the abstract sketch of pp. 39-40.
%
%   B. pp. 43-64  THE CHAPTER IS AUDITED FROM OUTSIDE, and one of its
%                 definitions does not survive. Lê Dũng Tráng's errata to
%                 §§20-21 (Jan. 1968); his own dossier on regular sequences
%                 (pp. 49-51); the typescript PS in which he concedes that the
%                 definition of « immersion régulière » at IV 16.9.2 IS NOT
%                 THE RIGHT ONE outside the noetherian case and that the
%                 serviceable notion is the Koszul one; and Deligne's three
%                 letters of October and December 1966, which supply what
%                 0_IV is missing and break the proof of 16.1.5.
%
%   C. pp. 66-98  COMPLÉMENTS FILED BY PARAGRAPH: « Divers » (Main Theorem,
%                 essentially of finite presentation, zero-dimensional
%                 preschemes, the condition (h_0), π_0 and étale coverings,
%                 two typescript leaves), then « IV 5 » (dimension, a
%                 ∂-category, finiteness of the integral closure) and
%                 « IV 12 » (a non-Cohen-Macaulay locus that jumps the wrong
%                 way, and a proposition on graded algebras).
%
% WHAT ONLY A FOLDER-WIDE READING CAN SAY. Four things, none of them visible
% inside any one batch:
%
%   1. THERE ARE THREE DELIGNE LETTERS, NOT FOUR. Pages 61-64 are not a fourth
%      letter: they are pages 3 to 6 of the letter of 17 October 1966 whose
%      pages 1 and 2 are pages 59-60. The join is exact — page 60 ends by
%      constructing A' = A[[T_i]] and page 61 opens « Remplaçant A par A' » —
%      and it is confirmed independently by page 64, where Deligne writes « la
%      démonstration tordue donnée pg 2 » of the commutation of Tor with the
%      pro-object completion: his page 2 is folder page 60, and that
%      proposition is exactly what stands there. Batches 3 and 4 each say the
%      letter continues into the other; neither can state the join.
%
%   2. THE LETTERS ARE FILED IN REVERSE DATE ORDER, and the reading restores
%      the chronology (17 Oct., 29 Oct., 15 Dec.), which means reading pages
%      59-64 before 57-58 before 55-56. Read in the folder's order the
%      correspondence is incoherent — the 15 December letter answers a letter
%      of his of the 28th, on a hypothesis the 29 October letter raised.
%
%   3. THE REGULAR-SEQUENCE QUESTION IS ASKED IN BATCH 3 AND ANSWERED IN
%      BATCH 4. Page 51 lists the paragraphs on regular sequences, points to
%      « Lettres Deligne 17. Oct. 66. et 29 Oct. 66. », and carries a boxed,
%      struck-out question: « Notion d'immersion régulière dans le cas non
%      noethérien : est-elle vraisemblable ?? Vérifier § 16 !! ». The
%      mathematics that answers it — the equivalence of the regular, Koszul
%      and Ext-vanishing conditions, with the Nagata prime-avoidance lemma —
%      is on pages 62-63, inside the letter of 17 October; and the verdict is
%      the typescript of pages 52-53. Grothendieck's own hand answers in the
%      margin of page 62: « garder EGA IV 19, à compléter éventuellement ».
%
%   4. THE BRACE SUPERSCRIPT IS ONE NOTATION ACROSS THE FOLDER. X_0^{{Λ}} is
%      introduced at the foot of page 21 and reused at pages 27, 28, 30, 31
%      and 32 for Y^{{Δ^n}}, X^{{S'}}, X_Y^{{Y_{S'}}}. It is one construction,
%      the Weil restriction, and it is the same construction the marginal of
%      page 20 attaches a proper name to: « Greenberg ».
%
% NOTATION fixed for the folder.
%   - The brace superscript X^{{Λ}} is kept as he writes it and glossed once,
%     at page 21, as the Weil restriction along Spec Λ → S.
%   - The order of a jet is written as a PARENTHESISED superscript throughout
%     — Δ^{(n)}, π^{(n)}, P^{(n)}, d^{(n)}, ε^{(n)}, φ^{(n,n')} — where the
%     pages alternate freely between (n) and n. The choice is ours and is
%     footnoted at page 26, where the transcription flags the alternation.
%   - His underlined sheaf letters — O_S, O_X, Λ, I, Γ, S, Aut, Is, Hom — are
%     set upright, without the underline.
%   - Aut^{X_0}_T(X), Is^{X_0}_T(ξ,X), S(X̄/X_0) are kept as he writes them.
%   - π(U) primitive elements, γ(U) group-like elements, π_0 and γ_0 their
%     kernels over U; δ is the REDUCED diagonal on I(U), which the pages use
%     and never define.
%   - L_k(M) is the Lie algebra of k-endomorphisms of M, as pages 13-14 use
%     it.
%   - OF(S) is the set of open-and-closed subsets of S; the pages write it as
%     one ligatured sign.
%   - On pages 74 to 79 every statement carries a STRONGER VARIANT IN
%     PARENTHESES — « injectif (bijectif) », « surjectif (bijectif) »,
%     « fidèle (pleinement fidèle) ». Each list is therefore two lists, and is
%     read here twice. This key is his, is constant over six leaves, and is
%     what makes those leaves legible.
%   - On pages 86-87 the letters p and q change sides between a lemma's
%     statement and its proof. Fixed here once for all: 𝔮 is a prime of B and
%     𝔭 = 𝔮 ∩ A. Footnoted at page 86.
%   - prof, Codim, ht, Inf, Sup, rg, C.M., 0_IV and the EGA reference numbers
%     are kept as written.
%
% SUBSTANTIVE DEPARTURES, each footnoted where it occurs.
%   - p. 5: the degree of the discriminant is written « (k−1)(r+1) », and the
%     marginal brace announcing it computes (k−1)^{r+1}. Neither is right for
%     r ≥ 2; the degree is (r+1)(k−1)^r. The written form is correct for r = 1
%     and for no other r.
%   - p. 4: the empty parenthesis in « Ω^n_{X/S} ≃ f^*( – ) » is filled here;
%     the completion is ours and is the only one the preceding formula allows.
%   - p. 8: « log gg' = log g log g' » is written with no + on the leaf. The
%     true identity is log(gg') = log g + log g'.
%   - p. 12: « exp ∘ α : π(U) ⊗ I(Λ) → 1 + π_0(U ⊗ Λ) » — the target is
%     γ_0(U ⊗_k Λ), the group-like elements, not 1 + π_0.
%   - p. 15: « induisent 0 sur B ⊗_Λ k ≃ Λ » — B ⊗_Λ k is A, not Λ, since
%     B = A ⊗_k Λ. The transcription flags the reading.
%   - p. 24: « Γ_φ = φ ×_S id_X : X' → X' ×_S X » does not typecheck. The
%     graph of φ : X' → X is (id_{X'}, φ). Transcribed unrepaired, repaired
%     here.
%   - p. 34: the constant of the first counterexample is c in the body and a
%     in the margin, and the intermediate identification of A is a discarded
%     simplification. Stated here in the form the conclusion needs.
%   - p. 83, Corollaire 2: the x_i are said to be part of a system of
%     parameters « de A/𝔭 », where they lie in 𝔭 and are therefore zero in
%     A/𝔭. Read as A, and the completion needs dim A/𝔭 + dim A_𝔭 = dim A,
%     which is supplied.
%   - p. 83, Corollaire 1: the statement hesitates between a local ring and a
%     noetherian ring with an ideal 𝔪; the final clause « si A local » settles
%     it for the second reading, and it is taken that way.
%   - p. 72: the list was renumbered from (v…) to (iv…) and the implications
%     below it were not; one line of the chain, « (iv) ⟺ (v) », has no
%     referent in the final numbering and is reported as such.
%   - p. 97, Question (i): « A plat/K ⟹ C plat/B » does not follow the pattern
%     of (ii) and (iii). Left as written, flagged; it is a question, not an
%     assertion.
%
% WHAT THE FOLDER ANNOUNCES AND NEVER ESTABLISHES.
%   - The sorites on the exponential are numbered (1), (2), (3), (4) and then
%     (6): no (5) is in the folder.
%   - The two torsor theorems (pp. 20-21, p. 28) are never turned into a
%     cohomological classification; H^1 is not taken anywhere in the folder.
%   - Page 28 asks, of the arrow P^n(Y/S) → f^*P^n(X/S), « comment ? », and
%     does not answer.
%   - Page 48 breaks off on his own ellipsis: « Corollaire. Si f : X → Y
%     est … ».
%   - Pages 39-40, « Foncteurs différentiels d'ordre n », never reach the
%     definition they are headed with.
%   - Page 70's Proposition has a (i) and an empty (ii).
%   - Page 97's Question — three implications about flatness — is asked and
%     not answered.
%   - Deligne's (v) ⟹ (iii) on page 62 is deferred (« Va être modulée ») and
%     what follows proves a different proposition; his struck N.B. on page 63
%     conjectures (iv) ⟹ (iii) without the hypothesis and does not prove it;
%     and on page 64 he says outright that he cannot prove the vanishing he
%     needs (« mais j'en suis incapable »).
%   - Three documents the folder cites and does not contain: the letters of
%     Grothendieck the three Deligne letters answer; the letter of Serre
%     enclosed with page 64; and Quillen's letter on Ω_{B/A} cited on page 62.
%     The addressee of the typescript of pages 52-53 is never named.
%
% LEGIBILITY, and where a human checker should go first. Pages 9 to 13 are a
% fast struck draft whose formulas carry and whose connective prose does not:
% the reading of pages 8-16 given here is a RECONSTRUCTION from the surviving
% formulas, and says so. Page 48 is the hardest leaf of the third batch. Pages
% 74-76 carry a palimpsest and a redrawn sketch; 77 and 79 are pencil over a
% typescript that shows through, and carry the most illegibles in the folder.
% Page 40 has three overwritten layers under two diagonals. The arrow
% directions of the lower diagram of page 27 are not on the leaf.
\documentclass[11pt,a4paper]{article}
% Le préambule commun à toutes les transcriptions du fonds Grothendieck.
%
% Il tient en une idée : **l'incertitude se compose, elle ne se résout pas.**
% Une transcription qui lisse un mot douteux a détruit la seule chose pour
% laquelle on la faisait — un lecteur doit pouvoir savoir, à chaque endroit,
% ce qui était sur le feuillet et ce qui a été ajouté. Les six macros qui
% suivent sont donc l'appareil critique, et non de la décoration.
%
% Les mêmes commandes sont reconnues par `scripts/render.mjs`, qui produit la
% vue de lecture du site. Une macro ajoutée ici doit l'être aussi là-bas,
% faute de quoi le rendu échouera bruyamment — ce qui est voulu : un rendu
% silencieusement faux est pire qu'un rendu absent.

\ProvidesPackage{grothendieck}[2026/08/08 Transcriptions du fonds Grothendieck]

\usepackage{amsmath,amssymb,amsthm}
\usepackage{mathrsfs}
\usepackage{stmaryrd}
% Les diagrammes commutatifs sont partout dans ce fonds : c'est la langue
% même de la géométrie algébrique de Grothendieck, et une transcription qui
% les rendrait en prose ne transcrirait rien.
\usepackage{tikz-cd}
% Français par défaut — c'est la langue des feuillets — mais l'anglais est
% chargé aussi : la traduction et le résumé sont des documents anglais, et la
% typographie française (espaces avant les deux-points) y serait une faute.
% Ces documents basculent par \selectlanguage{english} en tête de corps.
\usepackage[english,french]{babel}
\usepackage{fontspec}
\usepackage{geometry}
\geometry{a4paper,margin=25mm,marginparwidth=28mm}
\usepackage{marginnote}
\usepackage{xcolor}
% ulem, not soul. `soul`'s \st and \ul are built on a character-by-character
% scan that XeLaTeX's font machinery breaks: the document fails with an
% undefined control sequence rather than degrading. ulem does the same two jobs
% and is Unicode-engine safe. `normalem` stops it hijacking \emph, which the
% transcriptions use for Grothendieck's own emphasis.
\usepackage[normalem]{ulem}
\usepackage{hyperref}
\hypersetup{colorlinks=true,linkcolor=black,urlcolor=black,pdfborder={0 0 0}}

% --- Métadonnées du lot -----------------------------------------------------
% Reprises telles quelles par le script de rendu, qui les affiche en tête de
% la vue de lecture. Elles fixent ce que le fichier prétend couvrir : sans
% elles, une transcription flotte sans point d'attache dans un fonds de seize
% mille feuillets.

\newcommand{\folder}[1]{\gdef\grothFolder{#1}}
\newcommand{\batch}[1]{\gdef\grothBatch{#1}}
\newcommand{\leaves}[2]{\gdef\grothFirst{#1}\gdef\grothLast{#2}}
\newcommand{\foldertitle}[1]{\gdef\grothFoldertitle{#1}}
\gdef\grothFolder{?}\gdef\grothBatch{?}\gdef\grothFirst{?}\gdef\grothLast{?}\gdef\grothFoldertitle{}

% --- Le résumé --------------------------------------------------------------
%
% Il ouvre la lecture modernisée, sous le titre « Résumé », et il est composé
% en petit. Ce n'est pas un ornement : c'est le seul passage du document qui
% ne soit pas des mathématiques, et un lecteur doit pouvoir le reconnaître
% avant de l'avoir lu — puis le sauter s'il connaît déjà le sujet, ou s'y
% arrêter s'il ne le connaît pas. Le filet qui le suit marque la reprise du
% corps.

\newenvironment{resume}
  {\small\setlength{\parskip}{0.45em}}
  {\par\medskip\hrule\medskip}

% --- Mots-clés ---------------------------------------------------------------
%
% En anglais, en clôture du résumé de la lecture modernisée : le vocabulaire
% moderne sous lequel chercher ce que les feuillets construisent. Le manifeste
% du site (`npm run manifest`) les extrait pour étiqueter la cote entière —
% cette ligne est la source unique des tags, il n'y a pas de fichier à côté.
\newcommand{\keywords}[1]{\par\smallskip\noindent
  {\footnotesize\textsc{Keywords} --- \textit{#1}}\par}

% --- L'appareil critique ----------------------------------------------------

% Le numéro de feuillet, en marge. C'est l'ancre qui permet de régler un
% désaccord sur une formule en regardant *un* feuillet plutôt qu'une chemise
% de six cents. Le site s'en sert aussi pour tourner les pages du fac-similé
% quand on fait défiler la transcription.
\newcommand{\leaf}[1]{%
  \marginnote{\footnotesize\color{gray!70}\textsf{#1}}%
  \label{leaf:#1}\ignorespaces}

% Illisible : on ne devine pas. Un mot inventé qui se lit comme les autres est
% le pire résultat possible.
\newcommand{\ill}{\textcolor{red!60!black}{[\,\ldots\,]}}

% Lecture douteuse : le mot est donné, et signalé comme incertain.
\newcommand{\uncertain}[1]{\uline{#1}}

% Ajout de l'éditeur — une lettre manquante, un mot sous-entendu.
\newcommand{\add}[1]{\textcolor{blue!50!black}{[#1]}}

% Biffé par Grothendieck lui-même. Ce qu'il a rayé fait partie du document :
% une rature dit souvent où la pensée a changé de direction.
\newcommand{\struck}[1]{\sout{#1}}

% Note du transcripteur, sur l'état matériel du feuillet.
\newcommand{\note}[1]{\par\smallskip{\small\color{gray!60!black}\textit{[#1]}}\par\smallskip}

% Note marginale de Grothendieck, distincte du corps du texte.
\newcommand{\marginal}[1]{\par\smallskip\noindent\hspace{1em}%
  {\small\color{gray!60!black}#1}\par\smallskip}

% --- Titre ------------------------------------------------------------------

\AtBeginDocument{%
  \begin{center}
    {\large\bfseries Fonds Alexandre Grothendieck --- cote n\textsuperscript{o}~\grothFolder}\\[2pt]
    {\itshape\grothFoldertitle}\\[4pt]
    {\small feuillets \grothFirst--\grothLast\ (lot \grothBatch)}\\[2pt]
    {\footnotesize\color{gray!60!black}Transcription établie d'après le fac-similé numérisé
     par l'Université de Montpellier.\\
     Les lectures incertaines sont soulignées, les illisibles notées [\,\ldots\,],
     les ajouts entre crochets.}
  \end{center}
  \vspace{1em}}

\endinput


\folder{23}
\pages{1}{98}
\dating{1966-1968}
\watermark{Édition de démonstration\\interprétation personnelle de l'œuvre}
\foldertitle{EGA IV. Compléments : notes manuscrites (s.d.), lettres (1966, 1968), tapuscrit (s.d.). — lecture modernisée du dossier entier}

\begin{document}

\section*{Résumé}

\begin{resume}
Dériver, c'est comparer un objet avec lui-même déplacé d'une quantité si petite
qu'on convient de négliger son carré. Tout ce dossier part de là, et la première
chose qu'il fait est de prendre cette convention au mot : on décrète qu'il
existe une quantité $\varepsilon$ non nulle dont le carré est nul, et l'on
regarde ce que deviennent les objets de la géométrie quand on leur permet de
bouger de $\varepsilon$.

La réponse tient en une phrase, et c'est la phrase que ces feuillets écrivent
quatre fois dans quatre langages. \emph{Bouger d'un infiniment petit, c'est la
même chose que dériver.} Un automorphisme qui ne bouge rien au premier ordre
est une dérivation déguisée ; une dérivation, réciproquement, s'exponentie en un
automorphisme — exactement comme, en analyse, la donnée d'un champ de vitesses
équivaut à la donnée du mouvement qu'il engendre. La formule qui fait passer de
l'un à l'autre est la série de l'exponentielle, et elle ne demande rien d'autre
que de pouvoir diviser par les entiers.

Ce qui est plus surprenant, et qui occupe la moitié du dossier, est que cette
équivalence ne se contente pas d'être vraie : elle est \emph{rigide}. Quand on
déforme un objet géométrique d'un infiniment petit, les manières de le déformer
ne forment pas un ensemble quelconque. Elles forment ce qu'on appelle un
\emph{espace principal homogène} : un ensemble sur lequel un groupe agit de
telle sorte qu'on peut toujours passer d'un point à un autre, et d'une seule
façon. C'est la situation d'un plan sans origine — les vecteurs agissent sur les
points par translation, et dès qu'on choisit un point pour origine, le plan
\emph{devient} l'espace des vecteurs, mais le choix n'est pas donné avec le
plan. Les déformations d'un schéma sur une base infinitésimale sont dans ce
rapport-là avec les dérivations : il y a autant de déformations que de
dérivations, mais rien ne dit laquelle est la déformation nulle tant qu'on n'a
pas choisi une trivialisation. Le dossier démontre cela deux fois, pour les
déformations d'un schéma d'abord, puis pour un objet qu'il appelle une
\emph{connexion} — une manière de transporter, le long des directions de la
base, ce qui vit au-dessus d'elle.

Entre les deux, il construit l'outil : le \emph{jet}. Le jet d'ordre $n$ d'une
fonction en un point, c'est son développement de Taylor tronqué à l'ordre $n$ ;
et l'observation qui organise tout est qu'on peut fabriquer un espace
géométrique dont les points \emph{sont} les développements de Taylor. Il
s'obtient en épaississant la diagonale — l'ensemble des couples de points égaux
— jusqu'à l'ordre $n$ : un couple de points « égaux à l'ordre $n$ près » est
exactement ce qu'il faut pour parler de la variation d'une fonction entre eux.
Toute la théorie des opérateurs différentiels tient dans cette idée, et ces
feuillets en donnent la version où l'on ne suppose plus que les objets sont des
fonctions.

L'autre moitié du dossier ne construit rien : elle corrige. Et c'est là qu'il
faut dire ce que ce dossier est. Les \emph{Éléments de géométrie algébrique}
étaient, au moment de ces feuillets, un ouvrage en cours d'impression, et le
chapitre IV en était la partie la plus lourde. Ce que la chemise contient sous
le mot « compléments » est le courrier de ce chantier : les errata qu'un jeune
lecteur, Lê Dũng Tráng, relève en janvier 1968 dans les définitions des
paragraphes 20 et 21, en s'excusant de leur trivialité ; et trois lettres de
Pierre Deligne, de l'automne 1966, qui font tout autre chose que relever des
coquilles. Deligne lit le chapitre et renvoie ce qui casse. Une hypothèse de
séparation lui « semble idiote » et il en donne un contre-exemple. Une
démonstration est « canulée », et il dit en trois lignes pourquoi. Il comble au
passage ce qui manque : le théorème d'Auslander qui relie la dimension
projective d'un module à sa profondeur, un énoncé de relèvement d'un corps
résiduel inséparable, l'équivalence entre platitude et projectivité formelle. Il
avoue aussi, deux fois, ce qu'il ne sait pas faire.

De cet échange sort une décision dont ces feuillets sont, à notre connaissance,
le dossier : la notion d'\emph{immersion régulière} telle que le chapitre IV la
définissait est mauvaise. Elle repose sur celle de \emph{suite régulière}, qui a
un défaut que Grothendieck énonce en une phrase : l'ordre des éléments compte,
et il est facile de donner deux éléments $x$, $y$ d'un anneau local tels que
$(x,y)$ soit une suite régulière et $(y,x)$ non. Une notion qui dépend d'un ordre
arbitraire ne se comporte bien ni par localisation, ni par descente, et une
définition bâtie dessus ne survit pas hors du cas noethérien. La bonne notion,
qui ne dépend d'aucun ordre, s'exprime par l'annulation d'un complexe — le
complexe de Koszul. Le dossier contient la lettre qui apporte les théorèmes
nécessaires, la fiche de travail où la question est posée puis rayée, et le
post-scriptum dactylographié où il tranche, propose d'ajouter aux épreuves un
aveu en bonne et due forme, et va jusqu'à proposer que les deux notions
échangent leurs noms — « il s'imposera alors de réserver le mot le plus court à
la notion qui manifestement est la plus serviable ». C'est ce qui est arrivé.

Il reste, pour finir, ce que ces feuillets montrent et qu'aucun livre ne montre :
un texte mathématique encore vivant. Deligne écrit qu'il a été malade et qu'il
lit le chapitre IV « comme lecture de convalescence ». Il avoue « avec honte »
qu'il serait prêt, « dans le même geste sacrilège », à supposer tous les schémas
séparés et noethériens. Il demande, à propos d'un autre chapitre, s'il ne
faudrait pas « parler de topologie avant de parler d'espace topologique ». Et
Grothendieck répond dans les marges, au crayon, par un mot chaque fois :
« oui », « Err », « dans réédition du Chap I ? », « faire errata ». On voit, sur
ces pages, la différence entre un résultat et le travail qui le précède : le
résultat est ce qui reste quand on a cessé de se tromper, et ces feuillets sont
le moment d'avant.

Les noms sous lesquels on cherchera la suite : dérivations et automorphismes
infinitésimaux, restriction de Weil et foncteur de Greenberg, jets et parties
principales, connexions, déformations et torseurs, suites Koszul-régulières et
complexe de Koszul, formule d'Auslander-Buchsbaum, critère de Serre, et — pour
la seule note du dossier qui regarde ailleurs — les théorèmes de type Lefschetz
pour le groupe de Picard.

\keywords{adjunction formula, complete intersection, projective bundle, Euler
sequence, conormal sheaf, discriminant of a hypersurface, exponential map,
topologically nilpotent element, bialgebra, primitive element, group-like
element, Baker-Campbell-Hausdorff formula, Cartier-Milnor-Moore theorem,
universal enveloping algebra, Poincare-Birkhoff-Witt theorem, derivation,
infinitesimal automorphism, tangent sheaf, infinitesimal deformation,
square-zero extension, Weil restriction, Greenberg functor, principal
homogeneous space, torsor, jet bundle, sheaf of principal parts, infinitesimal
neighbourhood of the diagonal, connection, differential operator, geometrically
integral, geometrically regular, separable field extension, p-basis, imperfect
field, fibred category, effective descent morphism, etale sheaf, constructible
sheaf, regular sequence, quasi-regular sequence, Koszul complex,
Koszul-regular sequence, regular immersion, prime avoidance lemma,
Auslander-Buchsbaum formula, projective dimension, depth, Tor-dimension,
pro-object, Artin-Rees lemma, formally projective module, local flatness
criterion, purely inseparable extension, Zariski main theorem, quasi-finite
morphism, essentially of finite presentation, von Neumann regular ring,
constructible topology, Serre's criterion, finiteness of integral closure,
Cohen-Macaulay locus, Poincare series, graded algebra, translation functor,
Lefschetz condition for the Picard group}
\end{resume}

\pagerange{1}{98}
\section*{Le fil du dossier, et les conventions}

\emph{Ce que la chemise contient.} Quatre-vingt-dix-huit feuillets, dont onze
ne portent qu'un titre de sa main en haut à droite et servent de chemises
intérieures. Ces onze titres sont la table des matières du dossier, et il vaut
mieux les avoir sous les yeux avant d'entrer :

\begin{enumerate}
  \item[\emph{p.\,1}] « EGA IV, Calcul diff. (compléments) » — couverture de toute la
  cote.
  \item[\emph{p.\,2}] « Discriminant etc » — couvre les pages 3 à 5.
  \item[\emph{p.\,7}] « Exponentielle » — couvre les pages 8 à 16.
  \item[\emph{p.\,17}] « Fibrés des germes d'ordre $n$ de sections, etc. Transformations
  infinitésimales » — couvre les pages 18 à 22, et, de fait, la suite jusqu'à
  la page 32.
  \item[\emph{p.\,33}] « Exemples, cas particuliers » — couvre les pages 34 à 36.
  \item[\emph{p.\,38}] « Foncteurs différentiels d'ordre $n$ » — couvre les pages 39 et
  40.
  \item[\emph{p.\,49}] « Suites régulières / Th. 11.2.9 de Raynaud » — couvre les pages
  50 à 53.
  \item[\emph{p.\,54}] « Lettres Deligne » — couvre les pages 55 à 64.
  \item[\emph{p.\,65}] « Divers …. » — couvre les pages 66 à 80.
  \item[\emph{p.\,81}] « IV 5 », puis « Compléments d'ordre divers » — couvre les pages
  82 à 89.
  \item[\emph{p.\,90}] « IV 12 », puis « Fibres des morphismes plats de présentation
  finie » — couvre les pages 91 à 98.
\end{enumerate}

Les errata de Lê Dũng Tráng (pages 43 à 46) et le feuillet « EGA IV 18 ; EGA
VI » (page 48) ne sont couverts par aucune chemise : ils sont glissés entre
celle des « Foncteurs différentiels » et celle des « Suites régulières », et
c'est le seul indice de classement qu'on ait sur eux. Les pages 6, 41, 42, 47,
68, 71, 85 et 94 sont blanches. La page 37 est un tapuscrit de rédaction
bourbachique, « n\textsuperscript{o} 215 », étranger au dossier : elle n'est
pas lue.

\emph{Les trois mouvements.} Le premier, des pages 3 à 40, est de sa main et
construit. Le deuxième, des pages 43 à 64, n'est pas de sa main — deux
correspondants et un tapuscrit — et corrige. Le troisième, des pages 66 à 98,
est de nouveau de sa main et range : des compléments classés par paragraphe du
chapitre.

Ce qui relie les deux premiers est le paragraphe 16 des \emph{Éléments}, qui
s'appelle « Calcul différentiel » et que la couverture du dossier nomme. Le
premier mouvement en construit l'appareil — les voisinages infinitésimaux de la
diagonale, les parties principales, les automorphismes infinitésimaux. Le
deuxième en démolit une définition, celle de 16.9.2. Le même paragraphe est
l'objet des deux, et c'est ce qui fait de cette chemise un dossier plutôt qu'un
tiroir.

\emph{La forme que le premier mouvement répète.} Il faut la dire une fois,
parce qu'elle revient quatre fois et qu'aucun feuillet ne la nomme. On se donne
une base infinitésimale — un anneau ou un schéma dont on a décrété qu'une
certaine puissance d'un idéal est nulle. Au-dessus d'elle, un objet, et sa
réduction quand on annule l'infiniment petit. Alors :

\begin{itemize}
  \item l'ensemble des \emph{manières de relever} l'objet réduit, ou de le
  trivialiser, ou de le transporter,
  \item est un \emph{espace principal homogène} sous le faisceau des
  \emph{automorphismes infinitésimaux} de l'objet trivial,
  \item et ce faisceau est un \emph{module de dérivations}.
\end{itemize}

Aux pages 8 à 16 l'objet est une algèbre et le troisième point est un théorème
(l'exponentielle) ; aux pages 18 à 22 l'objet est un schéma et c'est le
deuxième point qui est le théorème ; aux pages 23 à 29 l'objet est une connexion
et c'est encore le deuxième ; aux pages 30 à 32 l'objet est un morphisme et le
dossier n'énonce plus que la première moitié. Le troisième point n'est jamais
repris au niveau des schémas : le dossier ne calcule nulle part le faisceau
$\mathrm{Aut}^{X_0}_T(\xi)$ en termes de $\mathcal{T}_{X_0/S}$, et ne prend
nulle part la cohomologie qui transformerait ces théorèmes de torseur en une
classification. C'est la lacune principale du dossier, et elle est
délibérée : il construit l'objet et s'arrête avant de l'exploiter.

\emph{Les notations, valables pour tout le dossier.} Ses lettres soulignées —
$\mathcal{O}_S$, $\mathcal{O}_X$, $\Lambda$, $I$, $\Gamma$, $S$,
$\mathrm{Aut}$, $\mathrm{Is}$ — sont écrites ici sans soulignement, sauf
$\underline{\mathrm{Hom}}$ et $\underline{P}$, où le trait ne relève pas de son
habitude mais note le \emph{Hom interne}, c'est-à-dire le faisceau des
homomorphismes, et le faisceau des parties principales : c'est encore la
notation d'aujourd'hui, et l'effacer perdrait une distinction. L'exposant entre accolades $X^{\{\Lambda\}}$ est le sien et est
gardé : il est introduit au bas de la page 21 et désigne le schéma des points
de $X$ « de type $\Lambda$ », c'est-à-dire la restriction de Weil le long de
$\mathrm{Spec}\,\Lambda \to S$.\footnote{C'est le foncteur que M. Greenberg
avait introduit en 1961 dans le cas où $\Lambda$ est un anneau de vecteurs de
Witt tronqués, et le nom propre que la marge de la page 20 porte, lu
« Greenberg » sans certitude par le transcripteur, désigne selon toute
vraisemblance cette construction : elle est celle de la page même, et la marge
la commente. On le dit ici comme un appui à une lecture douteuse, non comme
une preuve.} L'ordre d'un jet est écrit partout en exposant \emph{entre
parenthèses}, $\Delta^{(n)}$, $\pi^{(n)}$, $P^{(n)}$, $d^{(n)}$,
$\varepsilon^{(n)}$, alors que les feuillets alternent librement entre $(n)$ et
$n$ selon le symbole ; le choix est nôtre et ne prétend qu'à la
constance.\footnote{La transcription signale l'alternance à la page 26 et ne la
régularise pas, ce qui est sa règle ; celle-ci est l'inverse, et il faut donc
savoir que la forme parenthésée sur $P$, $d$, $\varepsilon$ et $\varphi$ ne se
trouve pas sur les feuillets.} $\mathrm{OF}(S)$ désigne l'ensemble des parties
ouvertes et fermées de $S$, que les pages 74 à 79 écrivent d'un seul signe
ligaturé.

Une clef de lecture, enfin, qui vaut pour les pages 74 à 79 et sans laquelle
elles ne se lisent pas : chaque condition y est écrite sous sa forme faible,
avec la forme forte \emph{entre parenthèses} — « injectif (bijectif) »,
« surjectif (bijectif) », « fidèle (pleinement fidèle) ». Chacune de ces listes
est donc deux listes, et les équivalences valent deux fois, une fois avec
chaque lecture. La convention est constante sur six feuillets et elle est de
lui.

\emph{Ce que le dossier annonce et n'établit pas.} Il faut le dire d'emblée,
parce qu'un dossier de travail promet beaucoup et qu'un lecteur qui cherche un
de ces énoncés perdrait son temps à le chercher ici.

\begin{itemize}
  \item Les sorites sur l'exponentielle sont numérotés par lui $(1)$, $(2)$,
  $(3)$, $(4)$, puis $(6)$ : \emph{aucun $(5)$ n'est dans le dossier}.
  \item Les deux théorèmes de torseur, celui des pages 20 et 21 et celui de la
  page 28, ne sont jamais convertis en une classification cohomologique. Le
  dossier ne prend $H^{1}$ nulle part, et ne calcule nulle part le faisceau des
  automorphismes infinitésimaux en termes du faisceau tangent.
  \item La page 28 pose, d'un point d'interrogation et du mot « comment ? », la
  question de la flèche $P^{n}(Y/S) \to f^{*}P^{n}(X/S)$, et ne la résout pas.
  \item La page 48 s'interrompt sur ses propres points de suspension :
  « Corollaire. Si $f : X \to Y$ est … ».
  \item Les pages 39 et 40, sous le titre « Foncteurs différentiels d'ordre
  $n$ », n'atteignent jamais la définition qu'elles annoncent.
  \item La page 70 énonce une proposition dont le (i) est écrit et le (ii)
  resté vide.
  \item La page 84 pose l'axiome $\partial 1$ d'une $\partial$-catégorie ; le
  second axiome n'est pas dans le dossier.
  \item La page 97 pose trois implications sur la platitude et n'en démontre
  aucune.
  \item Dans la lettre du 17 octobre, l'implication (v) $\Rightarrow$ (iii) est
  ajournée (« Va être modulée ») et remplacée par un autre énoncé ; le N.B.
  qui la conjecturerait sans hypothèse de minimalité est biffé ; et Deligne dit
  lui-même, à sa dernière page, qu'il ne sait pas démontrer l'annulation dont il
  aurait besoin.
  \item Quatre documents sont cités et absents : les lettres de Grothendieck
  auxquelles les trois lettres de Deligne répondent ; la lettre de Serre jointe
  à la page 64 ; la lettre de Quillen sur $\Omega_{B/A}$ citée à la page 62 ; et
  la rédaction « copiée sur la primitive » que la lettre du 29 octobre dit
  joindre. Le destinataire du tapuscrit des pages 52 et 53 n'est pas nommé.
\end{itemize}

\pagerange{3}{5}
\section*{I. Le canonique d'une intersection complète, et le discriminant (pages 3 à 5)}

Trois feuillets sous la chemise « Discriminant etc ». Les deux premiers sont
écrits au net et se lisent d'un bout à l'autre ; le troisième est un brouillon.

\emph{La formule d'adjonction.} Soient $S$ un préschéma, $E$ un faisceau
localement libre de rang $r+1$ sur $S$, et $P = \mathbf{P}(E)$ le fibré
projectif associé, de morphisme structural $g : P \to S$.\footnote{Le morphisme
$g$ n'est nommé nulle part sur le feuillet ; les $g^{*}$ des formules le
supposent, et c'est la transcription qui le restitue.} On dispose de la suite
exacte d'Euler
\[
  0 \longrightarrow \Omega^{1}_{P/S} \longrightarrow g^{*}(E)(-1)
  \longrightarrow \mathcal{O}_{P} \longrightarrow 0 ,
\]
d'où, en prenant les déterminants,
\[
  \Omega^{r}_{P/S} \xrightarrow{\;\sim\;}
  g^{*}\bigl(\Lambda^{r+1}_{\mathcal{O}_S} E\bigr)(-r-1) .
\tag{1}
\]

Soit maintenant $X \subset P$ un sous-préschéma lisse sur $S$, de dimension
relative $n = r - p$, d'idéal $J$, et $f : X \to S$ son morphisme structural. La
suite exacte conormale
\[
  0 \longrightarrow J/J^{2} \longrightarrow \Omega^{1}_{P/S} \otimes \mathcal{O}_X
  \longrightarrow \Omega^{1}_{X/S} \longrightarrow 0
\]
donne de même
\[
  \Omega^{r}_{P/S} \otimes \mathcal{O}_X \;\simeq\;
  \Omega^{n}_{X/S} \otimes \Lambda^{p}_{\mathcal{O}_X}(J/J^{2}) ,
\tag{2}
\]
et en composant (1) et (2)
\[
  f^{*}\bigl(\Lambda^{r+1}_{\mathcal{O}_S} E\bigr)(-r-1)
  \;\simeq\; \Omega^{n}_{X/S} \otimes \Lambda^{p}_{\mathcal{O}_X}(J/J^{2}) .
\]
C'est la formule d'adjonction, sous la forme $\omega_{X/S} \simeq
(\omega_{P/S} \otimes \det \mathcal{N}_{X/P})|_X$.

\emph{Le cas d'une intersection complète.} Supposons $X$ définie par un
sous-module gradué $M \subset \mathrm{Sym}^{*}(E)$, localement facteur direct et
localement libre de rang $p$.\footnote{Le feuillet ajoute, entre crochets, que
$M$ n'est déterminé que modulo le carré $I^{2}$ de l'idéal qu'il définit, et que
« sa structure est commune » — dernier mot de lecture incertaine. La remarque
est juste et n'est pas utilisée dans la suite. Le module dont $M$ est
sous-module n'est pas répété après le signe $\subset$ sur la page ;
$\mathrm{Sym}^{*}(E)$ est une restitution de l'éditeur de la transcription.}
Si $I_d$ désigne la composante de degré $d$ de l'idéal, on a des homomorphismes
$g^{*}(I_d)(-d) \to J$, d'où des homomorphismes surjectifs $\coprod
g^{*}(M_d)(-d) \to J$ induisant un isomorphisme
\[
  \coprod f^{*}(M_d)(-d) \xrightarrow{\;\sim\;} J/J^{2} ,
\]
le faisceau conormal de $X$ dans $P$.\footnote{Le feuillet écrit ce dernier
terme $\mathrm{Conormal}_{P/X}$, l'indice dans cet ordre ; c'est bien le
conormal de l'immersion $X \hookrightarrow P$, que l'on note aujourd'hui
$\mathcal{C}_{X/P}$ ou $\mathcal{N}^{\vee}_{X/P}$.} Notant $M_{d_i}$ les
composantes non nulles de $M$, de degré $d_i$ et de rang $\alpha_i$, avec
$p = \sum \alpha_i$, il vient
\[
  \Lambda^{p}_{\mathcal{O}_X}\, J/J^{2} \;\simeq\;
  f^{*}\Bigl(\bigotimes_{i} \Lambda^{\alpha_i} M_{d_i}\Bigr)
  \bigl(-\textstyle\sum \alpha_i d_i\bigr) ,
\]
et donc
\[
  \Omega^{n}_{X/S} \xrightarrow{\;\sim\;}
  f^{*}\Bigl(\bigl(\Lambda^{r+1}_{S} E\bigr) \otimes
  \bigotimes_{i} \Lambda^{\alpha_i} \check{M}_{d_i}\Bigr)
  \bigl(\textstyle\sum \alpha_i d_i - r - 1\bigr) .
\tag{3}
\]

Lorsque $\sum \alpha_i d_i = r+1$, la torsion disparaît et le canonique relatif
est image inverse d'un faisceau sur la base :
\[
  \Omega^{n}_{X/S} \xrightarrow{\;\sim\;}
  f^{*}\Bigl(\bigl(\Lambda^{r+1}_{S} E\bigr) \otimes
  \bigotimes_{i} \Lambda^{\alpha_i} \check{M}_{d_i}\Bigr) .
\tag{4}
\]
La parenthèse est laissée \emph{vide} sur le feuillet, avec un simple
tiret.\footnote{Le membre donné ici en (4) est nôtre. C'est le seul que (3)
autorise : il suffit d'y faire $\sum\alpha_i d_i - r - 1 = 0$. Aucune autre
lecture de la parenthèse vide n'est compatible avec la formule qui la précède
de trois lignes.}

\emph{L'exemple de la cubique plane.} Soit $E$ localement libre de rang $3$,
donc $r = 2$, avec un seul $M_{d_i}$, de degré $d_i = 3$ et de rang
$\alpha_i = 1$ : $X$ est une courbe elliptique lisse sur $S$, et l'on est dans
le cas $\sum \alpha_i d_i = 3 = r+1$. La formule (4) donne
\[
  \Omega^{1}_{X/S} \xrightarrow{\;\sim\;} f^{*}\bigl(\Lambda^{3} E \otimes
  \check{M}\bigr) .
\]
Si $E$ admet une filtration de facteurs $\mathcal{O}_S$, $L^{\otimes 2}$,
$L^{\otimes 3}$, où $L$ est inversible avec $f^{*}(L) \simeq
\mathcal{T}_{X/S}$,\footnote{Le symbole après « $f^{*}(L) \simeq$ » est une
lettre cursive à boucle ; elle est lue $\mathcal{T}$, le faisceau tangent, et la
ligne suivante du feuillet, $\Omega^{1}_{X/S} \simeq f^{*}(\check{L})$, confirme
la lecture. Le feuillet 15 explicitera la même lettre par
$\mathrm{Hom}_{B}(\Omega^{1}_{B/\Lambda}, B)$, ce qui la fixe pour tout le
dossier.} alors $\Lambda^{3} E \simeq L^{\otimes 5}$, et de
$\Omega^{1}_{X/S} \simeq f^{*}(\check{L})$ on tire $\check{L} \simeq
L^{\otimes 5} \otimes \check{M}$, c'est-à-dire
\[
  M \simeq L^{\otimes 6} .
\]
C'est le degré classique de l'équation de Weierstrass dans le fibré de Hodge :
la forme cubique qui définit la courbe vit à la sixième puissance du faisceau
inversible qui mesure son canonique.

\emph{Le discriminant.} Le troisième feuillet est un brouillon — des formules
alignées à gauche, un bloc de prose, une insertion en diagonale dans la marge —
et n'établit rien de complet. Ce qu'il pose est ceci. On note $\Phi^{k}(P^{r})$
l'espace des hypersurfaces de degré $k$ de $P^{r}$, et $\delta_{r,k}$ son
discriminant, défini au signe près comme une section de
$\mathcal{O}_{\Phi^{k}(P^{r})}(N)$, où $N$ est le degré du discriminant en les
coefficients de la forme, soit
\[
  N = (r+1)(k-1)^{r} .
\]
Ce degré n'est pas celui du feuillet.\footnote{La page écrit $(k-1)(r+1)$, ici
et dans les deux exposants suivants, et porte en marge une accolade sous
$\underbrace{k-1 \ \ k-1 \ \cdots \ k-1}$ avec la mention « $r+1$ facteurs »,
qui annonce donc $(k-1)^{r+1}$. Les deux écritures de la page ne coïncident
même pas entre elles. Le degré du discriminant d'une forme de degré $k$ en
$r+1$ variables est $(r+1)(k-1)^{r}$ : pour $r = 1$ il vaut $2(k-1)$, que la
forme écrite sur la page redonne exactement, et c'est sans doute pourquoi elle
est écrite ainsi ; dès $r = 2$ elle est fausse — pour les cubiques planes le
degré est $3 \cdot 2^{2} = 12$, et non $6$. Le feuillet est un brouillon et ne
revient pas dessus.}

Le reste du feuillet est juste et lisible. Une section $\varphi$ de
$\Phi^{k}(P^{r})$ au-dessus de $S$ — c'est-à-dire une famille d'hypersurfaces
de degré $k$ paramétrée par $S$ — définit un faisceau inversible $\check{M}$ sur
$S$, image inverse de $\mathcal{O}_{\Phi^{k}(P^{r})}(1)$, et le discriminant
s'y transporte en une section $\delta$ de $\check{M}^{\otimes N}$. L'énoncé
visé est alors :
\[
  P^{r}_{\varphi} \text{ est lisse sur } S
  \iff \delta \text{ est partout non nulle.}
\]
Et si $J = J_{\varphi}$ est l'idéal que $\varphi$ définit sur $P^{r}$, on a
$J \simeq g^{*}(M)(-k)$, provenant de l'homomorphisme $\mathcal{O}(-k) \otimes
g^{*}(M) \to \mathcal{O}_{P^{r}}$ obtenu en restreignant
$\mathcal{O}(-k) \otimes g^{*}\mathrm{Sym}^{k}(E) \to
\mathcal{O}_{P^{r}}$.\footnote{Le feuillet écrit $\mathrm{Symm}$, et cette
graphie est gardée dans la transcription sans être régularisée ; elle l'est
ici. Un signe noirci surmonte le $M$ de $J \simeq g^{*}(M)(-k)$, mais la
formule encadrée qui suit donne $M$ sans dual, et c'est la lecture retenue.}

\pagerange{8}{16}
\section*{II. L'exponentielle : une dérivation est un automorphisme infinitésimal (pages 8 à 16)}

Neuf feuillets sous la chemise « Exponentielle », intitulés par lui
« Sorites sur l'exponentielle » et numérotés $(1)$, $(2)$, $(3)$, $(4)$ — avec
un $(6)$ à la page 12 et aucun $(5)$ dans le dossier. Les feuillets 8, 14, 15 et
16 sont écrits au net ; les feuillets 9 à 13 sont un brouillon rapide, biffé et
surchargé, dont les formules portent et dont la prose de liaison, pour
l'essentiel, ne porte pas.

Il faut le dire avant d'entrer : \emph{ce qui suit des pages 9 à 13 est une
reconstitution}. Les énoncés donnés ici sont ceux que les formules survivantes
imposent, dans l'ordre où elles viennent ; les phrases qui les articulaient
sont perdues, et l'on ne peut pas garantir qu'il les enchaînait ainsi. Les
pages 8, 14, 15 et 16, elles, se lisent.

\subsection*{(1) L'exponentielle dans une algèbre topologique (page 8)}

Soit $U$ une algèbre associative unitaire sur $\mathbf{Q}$, linéairement
topologisée, séparée et complète, et $N(U)$ l'ensemble de ses éléments
topologiquement nilpotents.\footnote{Ces qualificatifs sont écrits en petit sous
le titre du feuillet, à droite, et raccordés par un trait au mot
« algèbre » de la ligne suivante ; ils qualifient $U$, et la transcription les
rattache ainsi. Si la topologie de $U$ est discrète, $N(U)$ est l'ensemble des
éléments nilpotents, ce que le feuillet note en passant.} Les deux séries
\[
  \exp \xi = \sum_{n \geqslant 0} \frac{1}{n!}\,\xi^{n} ,
  \qquad
  \log(1+\eta) = \sum_{n > 0} \frac{(-1)^{n+1}}{n}\,\eta^{n}
\]
convergent et définissent deux bijections réciproques
\[
  N(U) \;\underset{\log}{\overset{\exp}{\rightleftarrows}}\; 1_U + N(U) .
\]
Elles échangent addition et multiplication lorsque les éléments commutent :
\[
  \begin{cases}
    \exp(\xi + \xi') = \exp \xi \, \exp \xi' & \text{si } [\xi, \xi'] = 0 , \\[4pt]
    \log(g g') = \log g + \log g' & \text{si } g g' = g' g .
  \end{cases}
\]
La seconde identité est écrite sur le feuillet sans signe $+$ visible, sous la
forme « $\log g g' = \log g \; \log g'$ ».\footnote{C'est une omission, non une
autre identité : le second membre serait un produit, et l'équivalence des deux
conditions « si », que le feuillet marque par une double flèche verticale
$\Updownarrow$ entre elles, n'aurait plus de sens. Le membre juste est la somme,
et c'est lui qui est écrit ici.}

Deux compléments, que le feuillet porte entre crochets avec une flèche les
renvoyant après la ligne suivante. Si $I$ est un idéal à gauche ou à droite de
$U$, fermé, les bijections se restreignent en $N(U) \cap I \rightleftarrows 1_U
+ N(U) \cap I$. Et si $\varphi : U \to V$ est un homomorphisme d'algèbres
topologiques, alors $\varphi(N(U)) \subset N(V)$ et
\[
  \varphi(\exp \xi) = \exp \varphi(\xi) , \qquad
  \varphi(\log g) = \log \varphi(g) .
\]
En particulier toute sous-algèbre fermée est stable pour $\exp$ et
$\log$.\footnote{Le mot de remplacement écrit au-dessus d'« anneaux » biffé ne
se lit pas ; « algèbres » est une conjecture du transcripteur, et le sens le
demande — l'énoncé a besoin de la $\mathbf{Q}$-structure.}

\subsection*{(2) Éléments primitifs et éléments multiplicatifs (pages 9 à 11)}

On suppose maintenant $U$ munie d'une structure de bigèbre sur $k$ :
augmentation d'idéal $I(U)$, et diagonale
\[
  \Delta : U \longrightarrow U \,\widehat{\otimes}_{k}\, U
\]
qui soit un homomorphisme d'algèbres augmentées. On note $\delta$ la
\emph{diagonale réduite} sur $I(U)$, c'est-à-dire
$\delta(\xi) = \Delta(\xi) - \xi \otimes 1 - 1 \otimes \xi$.\footnote{Le
feuillet 9 écrit $\delta\xi = \xi \otimes \xi$ sous $\Delta(g) = g \otimes g$,
et n'a défini $\delta$ nulle part. La définition donnée ici est celle qui rend
la ligne vraie : si $g = 1 + \xi$, alors $\Delta(g) = g \otimes g$ s'écrit
$1\otimes 1 + \Delta(\xi) = 1 \otimes 1 + \xi\otimes 1 + 1\otimes\xi +
\xi\otimes\xi$, soit exactement $\delta(\xi) = \xi \otimes \xi$. C'est aussi
elle qui donne son sens à « primitif », $\delta(\xi) = 0$.}

On pose alors
\[
  \gamma(U) = \{\, g \in 1 + I(U) \;:\; \Delta(g) = g \otimes g \,\} ,
  \qquad
  \pi(U) = \{\, \xi \in I(U) \;:\; \delta(\xi) = 0 \,\} ,
\]
les éléments \emph{multiplicatifs} (aujourd'hui : de type groupe) et les
éléments \emph{primitifs}.\footnote{Trois notations sont écrites l'une sous
l'autre sur le feuillet 9 pour le premier ensemble : $\mathfrak{g}(U)$, biffé,
$\gamma(U)$, et une troisième biffée par-dessus un $\pi$ souligné. C'est
$\gamma(U)$ qui est repris aux feuillets 10 à 12, et c'est lui qui est écrit
ici. La lettre gothique $\mathfrak{g}$ servira, dans la marge du feuillet 12, à
tout autre chose : l'algèbre de Lie $\pi(U)$ elle-même.} Que $\gamma(U)$ soit
un \emph{groupe} est vérifié sur le feuillet par les deux calculs
\[
  \Delta(g g') = \Delta(g)\,\Delta(g') = (g \otimes g)(g' \otimes g')
  = g g' \otimes g g' , \qquad
  \Delta(g^{-1}) = \Delta(g)^{-1} = g^{-1} \otimes g^{-1} ,
\]
et la formation de $\gamma(U)$ est fonctorielle en le couple $(k, U)$. Que
$\pi(U)$ soit une \emph{algèbre de Lie} est immédiat sur la définition et n'est
pas vérifié sur le feuillet.

Si de plus $U$ est une $\mathbf{Q}$-algèbre, alors $\exp$ et $\log$ échangent
les deux :
\[
  \pi(U) \;\underset{\log}{\overset{\exp}{\rightleftarrows}}\; \gamma(U)
\]
sont des bijections réciproques. Il note lui-même ce qu'il suffit de vérifier :
que $\exp$ envoie $\pi(U)$ dans $\gamma(U)$ et $\log$ envoie $\gamma(U)$ dans
$\pi(U)$, le reste étant déjà acquis au point (1). On \emph{transporte} alors la
loi de groupe de $\gamma(U)$ sur $\pi(U)$, et la loi obtenue s'exprime en termes
de crochets par la formule universelle de Campbell-Hausdorff.\footnote{Le nom
est lu « Campbell-Hausdorff » dans le corps du feuillet 10 et
« Hausdorff-Campbell » dans sa marge, biffé, puis « Hausdorff-Campbell » aux
feuillets 14 et 15 ; l'ordre n'est pas fixé dans le dossier. La justification
que le feuillet esquisse entre crochets — par la bigèbre libre complétée sur un
ensemble — est celle d'aujourd'hui : la formule est universelle parce qu'elle
vaut dans l'algèbre de Lie libre complétée, où elle se calcule une fois pour
toutes.}

\subsection*{(3) Changement de base par une algèbre infinitésimale (pages 11 et 12)}

Soit $\Lambda$ une $k$-algèbre \emph{infinitésimale} : augmentée, d'idéal
d'augmentation $I(\Lambda)$ nilpotent. Alors $U \otimes_{k} I(\Lambda)$ est
topologiquement nilpotent dans $U \otimes_k \Lambda$, et le point (1) s'applique :
\[
  U \otimes_k I(\Lambda)
  \;\underset{\log}{\overset{\exp}{\rightleftarrows}}\;
  1 + U \otimes_k I(\Lambda) ,
\]
et de même avec $I(U) \otimes_k I(\Lambda)$.

On note $\pi_{0}(U \otimes_k \Lambda)$ et $\gamma_{0}(U \otimes_k \Lambda)$ les
éléments primitifs, respectivement multiplicatifs, de $U \otimes_k \Lambda$ qui
deviennent triviaux dans $U$ par l'augmentation $U \otimes_k \Lambda \to
U$.\footnote{L'indice $0$ n'est pas explicité sur les feuillets ; le sens donné
ici est celui que la prose fragmentaire du feuillet 11 indique — « les primitifs
de $U \otimes_k \Lambda$ qui [s'annulent] dans $U \otimes_k \Lambda \to U$, et
les éléments multiplicatifs de $U \otimes_k \Lambda$ qui donnent $1$ » — et le
seul qui rende vraies les deux formules qui suivent.} On obtient alors des
bijections réciproques
\[
  \pi_{0}(U \otimes_k \Lambda)
  \;\underset{\log}{\overset{\exp}{\rightleftarrows}}\;
  \gamma_{0}(U \otimes_{k} \Lambda) .
\]

Reste à comparer $\pi_0(U \otimes_k \Lambda)$ à $\pi(U) \otimes_k I(\Lambda)$.
Il y a toujours un homomorphisme canonique
\[
  \alpha : \pi(U) \otimes_{k} I(\Lambda) \longrightarrow
  \pi_{0}(U \otimes_{k} \Lambda) ,
\]
et \emph{si $\Lambda$ est plate sur $k$}, c'est un isomorphisme : car $\pi(U)$
est le noyau de $\delta : I(U) \to I(U) \otimes_k I(U)$, c'est-à-dire qu'on a
une suite exacte
\[
  0 \longrightarrow \pi(U) \longrightarrow I(U)
  \xrightarrow{\;\delta\;} I(U) \otimes_k I(U) ,
\]
et le foncteur $- \otimes_{k} I(\Lambda)$ la préserve. D'où, en composant avec
l'exponentielle, un isomorphisme
\[
  \pi(U) \otimes_{k} I(\Lambda) \xrightarrow[\;\sim\;]{\;\exp \circ\, \alpha\;}
  \gamma_{0}(U \otimes_{k} \Lambda) .
\]
Le feuillet écrit ce second membre « $1 + \pi_{0}(U \otimes_k \Lambda)$ », tel
quel.\footnote{C'est une erreur d'écriture, non une autre assertion : $1 +
\pi_0$ n'est pas un groupe, et l'exponentielle n'y arrive pas. L'image de
$\pi_0$ par $\exp$ est $\gamma_0$, par la bijection de la ligne précédente du
feuillet même, et c'est $\gamma_0$ qui est écrit ici.}

La marge du feuillet 12, écrite de bas en haut sur toute la hauteur de la page
et en deux ou trois colonnes serrées, esquisse ce qui manque : le théorème de
Cartier, $U \simeq U(\mathfrak{g})$ pour $\mathfrak{g} = \pi(U)$, et
Poincaré-Witt, $\mathfrak{g} \simeq \pi(U(\mathfrak{g}))$.\footnote{C'est
aujourd'hui le théorème de Cartier-Gabriel-Milnor-Moore : sur un corps de
caractéristique nulle, toute bigèbre cocommutative connexe est l'algèbre
enveloppante de son algèbre de Lie des éléments primitifs, et les deux foncteurs
$\mathfrak{g} \mapsto U(\mathfrak{g})$ et $U \mapsto \pi(U)$ sont quasi-inverses
l'un de l'autre. Cartier l'avait publié en 1962, Milnor et Moore en 1965 ; le
feuillet est donc contemporain de la seconde. La marge n'en donne que les deux
noms et les deux formules, tout le reste étant illisible : elle enregistre
qu'il sait où cela mène, non une démonstration.}

\subsection*{(4) Le théorème : dérivations et automorphismes infinitésimaux (pages 13 à 16)}

Le point (4) spécialise. Soit $M$ un $k$-module, $\mathfrak{g} = L_k(M)$
l'algèbre de Lie de ses endomorphismes $k$-linéaires, et $U$ son algèbre
enveloppante, de topologie discrète. Pour $\xi$ nilpotent et $g = \exp \xi$, le
feuillet 13 aligne trois équivalences :
\[
  \xi . a = 0 \iff g a = a ,
  \qquad
  M' \text{ stable par } \xi \iff M' \text{ stable par } g
  \quad (M' \subset M) ,
\]
et, si $M$ est de plus muni d'une multiplication $(x,y) \mapsto x . y$,
\[
  \xi(x . y) = (\xi x) y + x (\xi y)
  \iff
  g(x . y) = (g x)(g y) .
\]
Sous les deux membres de cette dernière il écrit, soulignés, les deux mots qui
sont tout l'énoncé : \emph{dérivations} à gauche, \emph{automorphismes} à
droite.\footnote{La lecture des quatre mots portés sous l'équivalence
(« dérivations », « automorphismes », « multiplicatives », « invariants ») est
donnée comme douteuse par la transcription. Les deux premiers sont ceux que la
formule impose ; les deux seconds correspondent aux deux premières
équivalences de la même accolade.}

\emph{Proposition (pages 14 et 15).} Soient $k \supset \mathbf{Q}$, $A$ une
$k$-algèbre commutative, et $\Lambda$ une $k$-algèbre infinitésimale. Alors les
$\Lambda$-automorphismes de $A \otimes_{k} \Lambda$ qui induisent l'identité
modulo $I(\Lambda)$ correspondent biunivoquement aux $k$-dérivations
\[
  D : A \longrightarrow A \otimes_{k} I(\Lambda) ,
\]
la loi de groupe du premier ensemble se transportant sur le second par la
formule de Hausdorff-Campbell.\footnote{L'énoncé du feuillet 15 est accolé d'un
trait vertical ondulé dans la marge gauche, qui se poursuit sur les trois
premières lignes du feuillet 16 : c'est le résultat que la suite de feuillets
visait. Le mot qui qualifie $\Lambda$ dans son énoncé est une abréviation de
trois ou quatre lettres qui ne se lit pas ; « infinitésimale » est ce que la
démonstration exige et ce que les feuillets 11 à 13 emploient partout.}

La vérification est le calcul du feuillet 15, et elle tient en une chaîne. Avec
$B = A \otimes_{k} \Lambda$ :
\[
  \mathcal{T}_{B/\Lambda}
  = \mathrm{Hom}_{B}(\Omega^{1}_{B/\Lambda}, B)
  \simeq \mathrm{Hom}_{B}(\Omega^{1}_{A/k} \otimes_{k} \Lambda,
    A \otimes_{k} \Lambda)
  \simeq \mathrm{Hom}_{A}(\Omega^{1}_{A/k}, A \otimes_{k} \Lambda)
  = \mathrm{Der}_k(A,\, A \otimes_{k} \Lambda) ,
\]
et les $\Lambda$-dérivations de $B$ qui induisent $0$ sur $B \otimes_{\Lambda} k
\simeq A$ s'identifient aux $k$-dérivations $A \to A \otimes_{k}
I(\Lambda)$.\footnote{Le feuillet écrit « $B \otimes_{\Lambda} k \simeq
\Lambda$ », et la transcription signale qu'on lit $\Lambda$ où l'on attendrait
$A$. C'est bien $A$ : $B \otimes_\Lambda k = (A \otimes_k \Lambda)
\otimes_\Lambda k = A \otimes_k k = A$. L'énoncé serait faux avec $\Lambda$, et
la suite immédiate du feuillet, qui fait porter les dérivations sur $A$, le
confirme.} Le reste est le point (4) : une dérivation s'exponentie en un
automorphisme, et réciproquement.

Le feuillet 16 en tire le cas où l'on veut un \emph{module} et non seulement un
ensemble : si $\Omega^{1}_{A/k}$ est libre de rang fini sur $A$, alors
\[
  \mathrm{Hom}_{A}(\Omega^{1}_{A/k},\, A \otimes_{k} I(\Lambda))
  \simeq \mathcal{T}_{A/k} \otimes_{k} I(\Lambda) ,
\]
de sorte que le groupe des automorphismes infinitésimaux de type $\Lambda$ de
$A$ est le $k$-module $\mathcal{T}_{A/k} \otimes_{k}
I(\Lambda)$.\footnote{Le second membre écrit sur le feuillet, en interligne et
relié par une accolade, est « $\pi_{0}(\mathcal{T} \otimes_{k} U)$ », de lecture
douteuse, et la ligne suivante, qui commence par un signe $=$, ne se lit pas. La
forme donnée ici est celle que les hypothèses du feuillet — $\Lambda$ libre de
rang fini sur $k$, $\Omega^{1}_{A/k}$ libre de rang fini sur $A$ — rendent vraie
et qui rend ces hypothèses nécessaires : sans finitude, le $\mathrm{Hom}$ ne se
laisse pas sortir du produit tensoriel. Le feuillet s'achève sur une ligne entre
crochets, « [alors $A/k$ est diff. simple] », dont les deux derniers mots sont
de lecture incertaine.}

Une remarque qui n'est pas du dossier et qu'il faut faire. L'hypothèse
$k \supset \mathbf{Q}$ n'est pas un confort : en caractéristique $p$ la série de
l'exponentielle n'a pas de sens, l'équivalence entre dérivations et
automorphismes infinitésimaux tombe, et ce qui la remplace — puissances divisées,
théorème de Cartier sur les groupes formels — est d'une autre nature. Les pages
34 à 36 du dossier, qui sont tout entières en caractéristique $p$, montrent
d'ailleurs le genre de phénomène que cette section n'atteint pas.

\pagerange{18}{22}
\section*{III. Déformations d'un schéma sur une base infinitésimale (pages 18 à 22)}

La chemise 17 porte, de sa main et souligné en partie, « Fibrés des germes
d'ordre $n$ de sections, etc. Transformations infinitésimales ». Elle couvre
les cinq feuillets qui suivent, et en fait toute la suite jusqu'à la page 32 :
les pages 21 et 22 achèvent la série de corollaires commencée page 20, et
c'est le premier des deux endroits où l'argument du dossier enjambe une limite
de lot.

\subsection*{Le cadre (page 18)}

Soient $S$ un préschéma et $\Lambda$ un faisceau localement libre de rang fini
de $\mathcal{O}_{S}$-algèbres augmentées, d'idéal d'augmentation $I$, tel que
$\mathrm{gr}^{I}(\Lambda)$ soit localement libre — donc $I^{n} = 0$ pour $n$
grand. Un tel $\Lambda$ est dit \emph{de type infinitésimal}. On pose
\[
  T = \mathrm{Spec}(\Lambda) ,
\]
schéma fini et plat sur $S$, muni d'une section $S \to T$, et
$\Lambda_{n} = \Lambda/I^{n+1}$, $T_{n} = \mathrm{Spec}(\Lambda_{n})$, de sorte
que
\[
  S \simeq T_{0} \subset T_{1} \subset T_{2} \subset \cdots \subset T_{N} = T .
\]
Soit $X$ un $T$-schéma plat et lisse sur $T$, $X_{n} = X \times_{T}
T_{n}$.\footnote{L'insertion au-dessus de « sur $T$ » est de deux mots, dont le
second seul se lit à peu près ; la transcription lit « plat simple ». « Simple »
était, dans la terminologie des premiers \emph{Éléments}, le mot d'alors pour ce
qu'on appelle aujourd'hui \emph{lisse}, et c'est ainsi qu'il est rendu ici. La
platitude est de toute façon impliquée par la lissité ; qu'il l'écrive quand
même est une habitude de ces années.} On considère enfin la \emph{déformation
triviale}
\[
  \xi = X_{0} \times_{S} T = X_{0} \otimes_{\mathcal{O}_S} \Lambda ,
  \qquad \xi_{n} = X_{0} \times_{S} T_{n} , \quad \xi_{0} = X_{0} .
\]

Topologiquement rien ne bouge : $|X| = |\xi| = |X_{0}| = |\xi_{0}|$. Tout se
passe dans le faisceau structural. $\mathcal{O}_{X}$ est un faisceau de
$f_{0}^{*}(\Lambda)$-algèbres sur $X_{0}$, augmenté vers
$\mathcal{O}_{X_0}$, et la platitude de $X$ sur $T$ se lit sur la filtration
$I$-adique de $\mathcal{O}_{X}$ : le gradué associé est localement libre sur
$\mathcal{O}_{X_0}$.

On note alors $\mathrm{Aut}^{X_0}_{T}(X)$ le faisceau sur $X_0$ qui, à un ouvert
$U$, associe le groupe des $T$-automorphismes de $U$ induisant l'identité sur
$X_{0} \cap U$.\footnote{Le feuillet écrit deux notations l'une sous l'autre
avec un tiret devant la seconde, $\mathrm{Aut}^{X_0}_{T}(X)$ et
$\mathrm{Aut}^{X_0}_{\Lambda}(X)$ : ce sont deux noms du même faisceau, selon
qu'on nomme la base par le schéma ou par l'algèbre. Le premier est employé
partout ici.} C'est aussi, la page le note, le faisceau des automorphismes de la
$f_{0}^{*}(\Lambda)$-Algèbre $\mathcal{O}_{X_0}$-augmentée $\mathcal{O}_{X}$ —
autrement dit, la version faisceautique exacte de ce que la section II vient de
calculer au niveau des algèbres.

\subsection*{Le schéma des sections, et le théorème (pages 19 et 20)}

On considère le $S$-schéma
\[
  \overline{X} = \underline{\mathrm{Hom}}_{T/S}(T, X) ,
\]
que la marge appelle « schéma des sections de $X$ sur $\Lambda$ » : c'est la
restriction de Weil de $X$ le long de $T \to S$. La section $S
\xrightarrow{\sim} T_{0} \to T$ définit
\[
  \underline{\mathrm{Hom}}_{T/S}(T, X) \longrightarrow
  \underline{\mathrm{Hom}}_{T_0/S}(T_{0}, X \times_{T} T_{0}) \simeq X_{0} ,
\]
d'où un morphisme de $S$-préschémas $\pi^{X}_{\Lambda} : \overline{X} \to
X_{0}$, compatible avec les restrictions aux ouverts : $\overline{U} \simeq
(\pi^{X}_{\Lambda})^{-1}(U)$. Par suite $\mathrm{Aut}^{X_0}_{T}(X)$ opère comme
faisceau d'automorphismes sur le $X_0$-schéma $\overline{X}$, donc aussi sur le
faisceau
\[
  \mathcal{S}(\overline{X}/X_{0})
\]
des sections de $\overline{X}$ au-dessus de $X_{0}$.

D'autre part on dispose du faisceau $\mathrm{Is}^{X_0}_{T}(\xi, X)$ des
isomorphismes de $\xi$ sur $X$ induisant l'identité sur $X_0$ : c'est un
faisceau principal homogène sous $\mathrm{Aut}^{X_0}_{T}(X)$ opérant à gauche et
sous $\mathrm{Aut}^{X_0}_{T}(\xi)$ opérant à droite. La déformation triviale
fournit une section privilégiée $s_{0}$ de
$\mathcal{S}(\overline{\xi}/\xi_{0})$, et l'on pose $u \mapsto u(s_{0})$.

\emph{Proposition (page 20).} L'homomorphisme
\[
  \mathrm{Is}^{X_0}_{T}(\xi, X) \longrightarrow
  \mathcal{S}(\overline{X}/X_{0}) , \qquad u \longmapsto u(s_{0}) ,
\]
est un isomorphisme de faisceaux, compatible avec l'opération à gauche de
$\mathrm{Aut}^{X_0}_{T}(\xi)$ et celle à droite de
$\mathrm{Aut}^{X_0}_{T}(X)$.\footnote{La seconde ligne de l'énoncé est en
partie biffée et récrite sur le feuillet, et plusieurs mots de la clause de
compatibilité ne se lisent pas ; ce qui est donné ici est ce que les deux
corollaires immédiats exigent. Le sens des deux opérations — laquelle à gauche,
laquelle à droite — est celui qu'impose la formule $u \mapsto u(s_0)$, où $u$
part de $\xi$ et arrive dans $X$.}

\emph{Corollaire 1.} $\mathcal{S}(\overline{X}/X_{0})$ est principal homogène
sous $\mathrm{Aut}^{X_0}_{T}(X)$, à droite.

\emph{Corollaire 2.} $\mathrm{Aut}^{X_0}_{T}(\xi)$ opère naturellement sur
$\mathcal{S}(\overline{X}/X_{0})$.

\emph{Corollaire 3 (page 21).} Il existe un homomorphisme naturel de faisceaux
\[
  \mathcal{S}(\overline{X}/X_{0}) \longrightarrow
  \mathrm{Is}_{X_{0}}(\overline{\xi}, \overline{X}) ,
\]
compatible avec les deux opérations, et il est \emph{injectif} — car
$\mathrm{Aut}^{X_0}_{T}(\xi)$ opère fidèlement dans $\overline{\xi}$.

C'est le premier des deux théorèmes de torseur du dossier, et il dit ceci : une
déformation de $X_0$ sur la base infinitésimale $T$, munie d'une trivialisation,
n'est pas autre chose qu'une section du schéma des sections ; les déformations
trivialisées forment un espace principal homogène sous les automorphismes
infinitésimaux, et il n'y a pas d'origine.

\subsection*{La lecture inverse, et la notation à accolade (pages 21 et 22)}

La fin du feuillet 21 retourne l'énoncé, en signalant lui-même que ce n'est pas
très utile : $\overline{X}/X_{0}$ peut être vu comme un schéma sur $X_0$ à
faisceau structural
\[
  \mathrm{Aut}^{X_0}_{T}(\xi) = G^{\Lambda}_{X_{0}} ,
\]
opérant sur le \emph{schéma-type} $\xi = X_{0} \times_{S} T$ ; et la donnée de
$X$ sur $T$ se réduisant suivant $X_0$ équivaut à la donnée d'un schéma
$\overline{X}$ sur $X_{0}$ muni d'un tel faisceau structural, le faisceau
principal associé n'étant autre que
$\mathcal{S}(\overline{X}/X_{0})$.\footnote{L'adjectif qui qualifie $X$ sur $T$
dans cette phrase est souligné d'un trait sur le feuillet et ne se laisse pas
décider : quatre ou cinq lettres, dont une haute à l'avant-dernière place. Il
s'agit selon toute vraisemblance de « plat », ou de la reprise de l'hypothèse
« plat simple » de la page 18, mais aucune reconstitution n'est proposée ici.}

C'est là qu'est introduite la notation qui servira dans tout le reste du
dossier : $\overline{\xi}$ est noté
\[
  X_{0}^{\{\Lambda\}} ,
\]
schéma des points de $X_0$ \emph{de type $\Lambda$}.\footnote{Le mot qui suit
« s'appelle » est illisible sur le feuillet, un mot biffé le précède, et les
deux dernières lignes de la page sont entamées par le bord du scan. « Schéma des
points de type $\Lambda$ » est ce que la transcription restitue de « des pts
propres de $X_0$ de type $\Lambda$ » ; l'adjectif « propres » n'a pas d'emploi
visible et pourrait être un lapsus. La construction, elle, n'est pas douteuse :
c'est $\underline{\mathrm{Hom}}_{S}(\mathrm{Spec}\,\Lambda, X_0)$.} Le feuillet
22, qui ne porte que quatre lignes et achève la phrase laissée en suspens au bas
de la page 21, donne les deux lectures d'un de ses points : \emph{soit} comme un
automorphisme de $X_{0} \otimes_{\mathcal{O}_S} \Lambda$ induisant l'identité
sur $X_{0}$, \emph{soit} comme la donnée d'un champ de sections.

\pagerange{23}{29}
\section*{IV. Jets, parties principales, connexions (pages 23 à 29)}

Aucune chemise ne sépare cette suite de la précédente : elle commence sans
transition au feuillet 23, sous un titre qu'il écrit en toutes lettres en tête
de la page — « éléments de sections d'ordre $n$ de $Y/S$ sur $X/S$ » — et qu'il
reprend, souligné, dans la marge du feuillet 24.

\subsection*{Le voisinage de la diagonale, et le schéma $\pi^{(n)}$ (page 23)}

$X \times_{S} X$ est muni de ses deux projections $\mathrm{pr}_{1}$,
$\mathrm{pr}_{2}$ et de l'immersion diagonale $\Delta_{X}$. On note
$\Delta^{(n)}_{X/S}$ le $n$-ième voisinage infinitésimal de cette immersion,
\emph{vu comme $X$-préschéma par $\mathrm{pr}_{1}$}, et
$i^{(n)}_{X} : \Delta^{(n)}_{X/S} \to X \times_{S} X$ l'immersion,
$\delta^{(n)}_{X/S} : \Delta^{(n)}_{X/S} \to X$ le morphisme structural. C'est
l'objet qui représente, géométriquement, la notion de « deux points égaux à
l'ordre $n$ près ».

Pour $Y$ un $S$-préschéma, on pose alors
\[
  \pi^{(n)}(Y/X/S) = \bigl(\mathrm{pr}_{2}\, i^{(n)}_{X}\bigr)^{*}(Y)
  = \Delta^{(n)}_{X/S} \times_{X \times_{S} X} (X \times_{S} Y) ,
\]
un $\Delta^{(n)}_{X/S}$-schéma.\footnote{Sur le feuillet, la notation
$P^{(n)}(Y/X/S)$ est d'abord écrite au-dessus de ce produit fibré et reliée à
lui par une accolade, puis biffée : c'est la notation qui sera réservée, à la
page 24, à l'image directe. La biffure est donc une correction et non une
hésitation.} Sa restriction le long de la diagonale redonne $Y$ :
\[
  X \times_{\Delta^{(n)}_{X/S}} \pi^{(n)}(Y/X/S) = Y ,
\]
d'où une section $\varepsilon^{(n)}(Y/X/S) : Y \to \pi^{(n)}(Y/X/S)$, et dans
l'autre sens le morphisme
\[
  d^{(n)}(Y/X/S) : \pi^{(n)}(Y/X/S) \longrightarrow X \times_{S} Y
  \longrightarrow Y .
\]
Ce dernier n'est \emph{pas} un $X$-morphisme en général, et la page le note :
il passe par $\mathrm{pr}_{2}$ alors que la structure de $X$-schéma vient de
$\mathrm{pr}_{1}$. C'est précisément l'écart entre les deux projections qui
porte toute l'information différentielle.\footnote{Les deux glyphes $d$ et
$\varepsilon$ sont presque identiques dans cette main, et la transcription note
que la lecture est décidée non par le trait mais par la relation
$\varepsilon^{(n)} d^{(n)} = \mathrm{id}$ de la page 25 et par le sens des
flèches des diagrammes. On suit ici cette lecture, qui est la seule cohérente.}

\subsection*{Le fibré des éléments de sections (pages 24 à 26)}

Supposons maintenant $\Delta^{(n)}_{X/S}$ \emph{localement libre sur
$X$}.\footnote{Le feuillet écrit d'abord « plat, de présentation finie », biffe,
et écrit « loc. libre sur $X$ ». C'est l'hypothèse qui permet l'image directe,
et elle est vérifiée dès que $X$ est lisse sur $S$ — auquel cas
$\Delta^{(n)}_{X/S}$ est le spectre du faisceau des parties principales
$\mathcal{P}^{n}_{X/S}$, localement libre de rang $\binom{n+\dim}{n}$.} On peut
alors former l'image directe
\[
  P^{(n)}(Y/X/S) = \bigl(\delta^{(n)}_{X/S}\bigr)_{*}
  \bigl(\pi^{(n)}(Y/X/S) \,/\, \Delta^{(n)}_{X/S}\bigr) ,
\]
que la marge du feuillet 24 nomme, sur trois lignes soulignées, « fibré des
éléments de sections d'ordre $n$ de $Y/S$ sur $X/S$ ». C'est encore une
restriction de Weil, cette fois le long de $\delta^{(n)}_{X/S}$, et elle est
caractérisée par l'adjonction : pour tout $X$-préschéma $\varphi : X' \to X$,
\[
  \mathrm{Hom}_{X}\bigl(X', P^{(n)}(Y/X/S)\bigr) \simeq
  \mathrm{Hom}_{\Delta^{(n)}_{X/S}}\bigl(X' \times_{X}
  \Delta^{(n)}_{X/S},\ \pi^{(n)}(Y/X/S)\bigr) ,
\]
et en particulier $\Gamma\bigl(\pi^{(n)}/\Delta^{(n)}_{X/S}\bigr) \simeq
\Gamma\bigl(P^{(n)}/X\bigr)$.

Un élément de cet ensemble, la page le décrit concrètement : on prend le graphe
\[
  \Gamma_{\varphi} = (\mathrm{id}_{X'},\, \varphi) :
  X' \longrightarrow X' \times_{S} X ,
\]
on en prend le $n$-ième voisinage infinitésimal $\Gamma_{\varphi}^{(n)} \to X'
\times_{S} X$, et l'on se donne une section de $X' \times_{S} Y$ au-dessus de
lui. Le feuillet écrit le graphe « $\Gamma_{\varphi} = \varphi \times_{S}
\mathrm{id}_{X}$ ».\footnote{Tel quel, $\varphi \times_S \mathrm{id}_X$ est un
morphisme $X' \times_S X \to X \times_S X$ et ne va pas de $X'$ dans $X'
\times_S X$ : la formule ne typecheck pas. Le graphe de $\varphi$ est
$(\mathrm{id}_{X'}, \varphi)$, et c'est ce qui est écrit ici. Le membre médian
du feuillet est récrit par-dessus un premier essai et la page n'est pas
corrigée ; la transcription l'a laissée telle quelle, ce qui est sa règle.}

Les $\Delta^{(n)}_{X/S}$ forment un système inductif de schémas sur $X$, de
limite $\Delta^{(\infty)}_{X/S}$, et l'on a pour $n' \geqslant n$
\[
  \pi^{(n)}(Y/X/S) = \pi^{(n')}(Y/X/S) \times_{\Delta^{(n')}_{X/S}}
  \Delta^{(n)}_{X/S} ,
\]
d'où un système transitif de $X$-morphismes
$\varphi^{(n,n')} : P^{(n')}(Y/X/S) \to P^{(n)}(Y/X/S)$ vérifiant
\[
  \begin{cases}
    \varepsilon^{(n)} \circ \varphi^{(n,n')} = \varepsilon^{(n')} , \\[3pt]
    \varphi^{(n,n')} \circ d^{(n')} = d^{(n)} .
  \end{cases}
\]

\subsection*{Les deux opérations sur les sections (page 25)}

Sur les faisceaux de sections, la construction donne deux morphismes en sens
inverse :
\[
  d^{(n)}(Y/X/S) : \Gamma(Y/X) \longrightarrow
  \Gamma\bigl(P^{(n)}(Y/X/S)/X\bigr) ,
\]
\[
  \varepsilon^{(n)}(Y/X/S) : \Gamma\bigl(P^{(n)}(Y/X/S)/X\bigr)
  \longrightarrow \Gamma(Y/X) ,
\]
avec
\[
  \varepsilon^{(n)} \, d^{(n)} = \mathrm{id} .
\]
C'est exactement le couple « jet d'ordre $n$ » et « valeur » : $d^{(n)}$ associe
à une section son développement tronqué, $\varepsilon^{(n)}$ oublie le
développement et ne garde que la section. Le feuillet construit $d^{(n)}$ en
passant par
\[
  \mathrm{Hom}_{X}\bigl(\Delta^{(n)}_{X/S}, Y\bigr) =
  \Gamma\bigl(Y \times_{X} \Delta^{(n)}_{X/S} \,/\, \Delta^{(n)}_{X/S}\bigr) ,
\]
et en vérifiant que la construction se recolle sur les sections locales.

\subsection*{Les connexions d'ordre $n$, et le second théorème de torseur (pages 27 et 28)}

Le feuillet 27 reprend la construction en termes d'anneaux, sur un diagramme
affine : $A \to A \otimes A$ muni de ses deux structures de $A$-algèbre, $B$
une $A$-algèbre, et $A'$ le quotient de $A \otimes A$ qui définit
$\Delta^{(n)}$.\footnote{Plusieurs flèches du diagramme inférieur de cette page
sont \emph{sans pointe} sur le feuillet, et la transcription les rend avec
pointe faute de pouvoir décider leur sens ; ce sont donc des directions
supplées, non lues. Le point est à vérifier sur le feuillet, et c'est l'un des
endroits du dossier où un lecteur ayant l'original sous les yeux gagnera le
plus.} La flèche en pointillé du diagramme supérieur porte l'exposant à
accolade $Y^{\{\Delta^{(n)}_{X/S}\}}$, de la même famille que le
$X_{0}^{\{\Lambda\}}$ de la page 21 : c'est la même construction, appliquée à la
base infinitésimale $\Delta^{(n)}_{X/S}$ au lieu de $\mathrm{Spec}\,\Lambda$.

Le feuillet 28 donne alors les deux lectures qui sont l'aboutissement de la
suite :
le groupe
\[
  \Gamma\bigl(Y^{\{\Delta^{(n)}_{X/S}\}}/Y\bigr)
\]
est celui des automorphismes infinitésimaux d'ordre $n$ des fibres de $Y/X$ ;
et le groupe
\[
  \Gamma\bigl(P^{(n)}(Y/X/S)/Y\bigr)
\]
est celui des automorphismes de $\pi^{(n)}(Y/X/S)$ sur
$\Delta^{(n)}_{X/S}$ induisant l'identité sur $Y$, c'est-à-dire l'ensemble des
\emph{connexions d'ordre $n$} de $Y/X/S$ — cette dernière égalité soulignée
d'un trait gras sur la page : c'est la définition qu'elle visait. Et l'énoncé :

\emph{Si $Y/X$ est lisse, les connexions d'ordre $n$ de $Y/X/S$ forment un
espace principal homogène sous le groupe des automorphismes infinitésimaux
d'ordre $n$ des fibres.}\footnote{Les trois lignes du feuillet qui portent cet
énoncé sont très lacunaires : les deux objets se lisent, « principal homogène »
se lit, et les mots qui les articulent ne se laissent pas tous lire. L'énoncé
donné ici est celui que ces fragments composent, et c'est aussi la forme qu'a
prise, dans la littérature, l'énoncé correspondant : une connexion n'existe pas
toujours, mais dès qu'il en existe une, l'ensemble des connexions est un espace
affine sous le module des automorphismes infinitésimaux — ici, sous
$\mathcal{T}_{Y/X} \otimes \Omega^{1}_{X/S}$ à l'ordre un. Le dossier ne fait
pas ce dernier calcul.}

C'est la même forme qu'aux pages 20 et 21, un cran plus haut. La page pose
ensuite la comparaison des parties principales relatives et absolues, et
s'arrête sur une question qu'elle n'élucide pas : de l'homomorphisme
d'$\mathcal{O}_{Y}$-algèbres augmentées $P^{n}(Y/X) \to f^{*}(P^{n}(X/S))$ elle
passe à une flèche $P^{n}(Y/S) \to f^{*}P^{n}(X/S)$ qui porte, sur le feuillet,
un point d'interrogation et, sous lui, de sa main, le mot « comment ? ».

Le feuillet 29 est une feuille de calcul entièrement barrée de quatre longues
diagonales, sur laquelle survivent la formule $\varphi(g\, d^{n}f \otimes m) = m
\otimes f\, d^{n}g$ et quelques produits tensoriels : c'est la traduction
algébrique du même passage, abandonnée.

\pagerange{30}{32}
\section*{V. Transformations infinitésimales d'un morphisme (pages 30 à 32)}

Le dernier des quatre registres, et le plus bref. Soit $f : X \to Y$ un
$S$-morphisme et $S'$ une base infinitésimale au-dessus de $S$. On appelle
\emph{transformation infinitésimale de type $S'$ de $f$} un couple $w = (u,v)$
de $S'$-automorphismes de $X_{S'}$ et de $Y_{S'}$, induisant l'identité sur $X$
et sur $Y$, et rendant commutatif

\begin{tikzcd}
  X_{S'} \arrow[r, "u"] \arrow[d, "f_{S'}"'] & X_{S'} \arrow[d, "f_{S'}"] \\
  Y_{S'} \arrow[r, "v"] & Y_{S'}
\end{tikzcd}

Ces couples forment un \emph{groupe}, contenu dans le produit des groupes des
transformations infinitésimales de type $S'$ de $X$ et de $Y$ pris séparément —
c'est-à-dire que la condition de compatibilité avec $f$ découpe un sous-groupe
dans le produit, et rien de plus.

Ce groupe opère sur les sections. Si $g$ est une section de $X/Y$, on pose
\[
  \widetilde{w}(g) = u \, g_{S'} \, v^{-1} ,
\]
qui est une section de $X_{S'}$ sur $Y_{S'}$ — la vérification est immédiate :
$f_{S'} u g_{S'} v^{-1} = v f_{S'} g_{S'} v^{-1} = v v^{-1} = \mathrm{id}$. Le
feuillet appelle $\widetilde{w}(g)$ le « produit » de $g$, et identifie
l'ensemble où il vit :
\[
  \Gamma(X_{S'}/Y_{S'}) \simeq
  \Gamma\bigl(X_{Y}^{\{Y_{S'}\}} \,/\, Y\bigr) ,
\]
retrouvant l'exposant à accolade de la page 21, cette fois relatif à $Y$.

Les pages 31 et 32 énoncent trois \emph{propriétés fonctorielles}, dont la
prose de liaison est très lacunaire et dont seules les formules portent.

\begin{enumerate}
  \item[1)] Le groupe des transformations infinitésimales de type $S'$ de
  $X/Y$ opère sur $\Gamma(X_{S'}/Y_{S'})$, et l'on a la flèche de
  spécialisation $\Gamma(X/Y) \to \Gamma(X_{S'}/Y_{S'}) \simeq
  \Gamma(X_{Y}^{\{Y_{S'}\}}/Y)$, l'ensemble d'arrivée étant muni de son groupe
  d'opérateurs.\footnote{L'accolade et les deux mots « ensemble : groupe
  d'opérateurs » sont écrits sous le membre médian, sur le feuillet ; le N.B.
  qui suit, souligné, n'est lisible que par fragments et paraît noter que la
  flèche n'est pas une bijection en général, seulement une flèche
  d'ensembles-à-opérateurs.}
  \item[2)] Un $Y$-morphisme $h : X \to X'$ compatible avec deux
  transformations $w = (u,v)$ et $w' = (u',v')$ — c'est-à-dire tel que $u'
  h_{S'} = h_{S'} u$ — vérifie
  \[
    h^{\{Y_{S'}\}}\bigl(\widetilde{w}(g)\bigr) = \widetilde{w}'(h \circ g)
    \qquad \text{pour toute } g \in \Gamma(X/Y) .
  \]
  \item[3)] Compatibilité avec les produits fibrés. Si $w$, $w'$ sont des
  transformations de $X/Y$ et $X'/Y$ au-dessus de la même $v$, elles en
  définissent une, $w''$, de $X \times_{Y} X'$, via
  $(X \times_{Y} X')_{S'} = X_{S'} \times_{Y_{S'}} X'_{S'}$, et l'on a
  \[
    \widetilde{w}''\bigl(g \times_{Y} g'\bigr) = \widetilde{w}(g) \times_{Y}
    \widetilde{w}'(g') ,
    \qquad
    \bigl(X \times_{Y} X'\bigr)_{Y}^{\{Y_{S'}\}} \simeq
    X_{Y}^{\{Y_{S'}\}} \times_{Y} X'^{\,\{Y_{S'}\}}_{Y} .
  \]
\end{enumerate}

Le feuillet 30 porte, dans un cadre dessiné à la main et traversé d'un trait
oblique, trois lignes dont il ne subsiste presque rien.\footnote{C'est,
avec le tiers supérieur de la page 40, l'endroit du dossier où un lecteur ayant
le feuillet sous les yeux gagnerait le plus : le cadre est manifestement une
reformulation de l'ensemble, et il est perdu. Rien n'en est reconstitué ici.} Et
le bas du feuillet 32 est une esquisse de trois ou quatre mots par ligne, sur
une loi de composition, qui ne se lit pas. Le registre des morphismes n'atteint
donc pas, dans le dossier, le théorème de torseur que les deux registres
précédents ont atteint : il en pose le cadre et les compatibilités, et s'arrête
là.

\pagerange{34}{36}
\section*{VI. Deux contre-exemples en caractéristique $p$, et une feuille de critères (pages 34 à 36)}

Sous la chemise « Exemples, cas particuliers », deux feuillets écrits au net —
les mieux lisibles du dossier avec les lettres — et une feuille de
récapitulation à l'italienne. Les deux contre-exemples sont calibrés, et il faut
dire d'abord sur quoi : ils séparent chacun une propriété d'une propriété
\emph{géométrique}, c'est-à-dire d'une propriété qui doit survivre à
l'extension du corps de base. C'est la matière des paragraphes 4, 6 et 17 du
chapitre IV.

\subsection*{(i) Intègre, avec $k$ algébriquement fermé dans le corps des fractions, et pourtant pas géométriquement intègre (page 34)}

Soient $k$ un corps imparfait de caractéristique $p$, et $a$, $b$, $c$ des
éléments de $k$. On pose
\[
  F = a X^{p} + b Y^{p} - c , \qquad A = k[X,Y]/(F) .
\]
Alors $A$ est un anneau \emph{régulier}, $k$ est \emph{algébriquement fermé}
dans son corps des fractions, et pourtant $A \otimes_{k} \overline{k}$
\emph{n'est pas intègre} : en effet, sur $\overline{k}$ on peut extraire les
racines $p$-ièmes, et en posant $Z = a^{1/p} X + b^{1/p} Y$ puis $T = Z -
c^{1/p}$, il vient $F = Z^{p} - c = T^{p}$, de sorte que
\[
  A \otimes_{k} \overline{k} \;\simeq\; \overline{k}[X][T]/(T^{p}) ,
\]
qui n'est même pas réduit.\footnote{Trois remarques sur ce feuillet, qui est
écrit au net mais dont la marge hésite. D'abord la constante : elle est écrite
$c$ dans le corps du feuillet et $a$ dans la marge, et la page ne les concilie
pas ; le calcul ne dépend pas du choix et $c$ est retenu, comme dans le corps.
Ensuite la note marginale « faire $Z = X + Y$, donc $F = Z^{p} - [\ldots]$ »,
surmontée de « $a X^{p} + b Y^{p}$ » et du mot « non » : la substitution
$Z = X+Y$ ne donne $Z^p$ que si $a = b$, et le « non » est vraisemblablement son
propre refus de cette simplification. Enfin la ligne « $A \simeq
k[Z]/(Z^{p}-c)[X]$ », écrite juste après la définition de $A$ : elle décrit le
cas dégénéré $b = 0$, où $A$ est le polynôme en une variable sur le corps
$k(c^{1/p})$ — cas dans lequel $k$ n'est précisément \emph{pas} algébriquement
fermé dans le corps des fractions, ce qui ruinerait le contre-exemple. Elle
n'est pas conciliée avec le reste du feuillet et n'est pas utilisée ici. Dans
la dernière formule il écrit $k$ et non $\overline{k}$, tel quel.}

L'enjeu est net une fois dit dans le vocabulaire d'aujourd'hui. Que $k$ soit
algébriquement fermé dans le corps des fractions assure que $A \otimes_k
\overline{k}$ est \emph{irréductible} ; cela n'assure pas qu'il soit
\emph{réduit}, ce qui demanderait que l'extension $\mathrm{Frac}(A)/k$ soit
séparable. Géométriquement intègre est strictement plus fort que intègre plus
« $k$ algébriquement fermé dedans », et cet exemple est l'écart.

\subsection*{(ii) Régulier, à corps des fractions séparable, et pourtant pas géométriquement régulier (pages 34 et 35)}

Soient $p \neq 2$, $a \in k$ non puissance $p$-ième, et
\[
  F = Y^{2} - X^{p} + a , \qquad A = k[X,Y]/(F) .
\]
Ici $\mathrm{Frac}(A)$ \emph{est} séparable sur $k$ : $y$ est de degré $2$ sur
$k(x)$, et $2$ est inversible.\footnote{Le feuillet ne porte de cette phrase que des bribes : « Cependant,
pour $p \neq 2$ », puis « irréductible », puis « Cela prouve », puis « une
extension séparable » ; l'argument donné ici — $K = k(x)(y)$
avec $y^2 = x^p - a$, séparablement engendrée sur $k$ dès que $p \neq 2$ — est
celui que cette phrase désigne.} Et $A$ est régulier. Le seul point qui puisse
être singulier est celui où $\partial F/\partial Y = 2Y$ s'annule, c'est-à-dire
l'idéal maximal
\[
  \mathfrak{m}' = (Y,\ X^{p} - a) = (Y,\ Y^{2} - X^{p} + a) ,
\]
de corps résiduel $k[X]/(X^{p}-a) = k(a^{1/p})$, purement inséparable de degré
$p$ sur $k$. Le calcul du feuillet 35 tranche : le carré $\mathfrak{m}''$ est
engendré par $Y^{2}$, $Y(X^{p}-a)$, $(X^{p}-a)^{2}$ et $F$, donc contient
$X^{p}-a = Y^{2} - F$, de sorte que
\[
  \mathfrak{m}'/\mathfrak{m}'^{2} \text{ est engendré par } Y \text{ seul} .
\]
Il est donc de dimension $1$ sur le corps résiduel, égale à
$\dim A_{\mathfrak{m}'}$ : le point est régulier.

Mais $A \otimes_{k} \overline{k}$ ne l'est pas. Sur $\overline{k}$, posant
$\lambda = a^{1/p}$ et $Z = X - \lambda$, on a $F = Y^{2} - Z^{p}$, et à
l'origine $Y^{2} - Z^{p} \in \mathfrak{m}^{2}$ puisque $p \geqslant 3$ : le
cotangent est de dimension $2$ pour une dimension $1$, le point est singulier.
Le feuillet conclut en deux mots, qui sont les bons : le point est
\emph{« simple, mais non absolument simple »} — régulier, mais non lisse. Et il
ajoute, comme dernière ligne du feuillet 35, ce qui reste à vérifier pour que
l'exemple soit complet : que $k$ soit algébriquement fermé dans
$A$.\footnote{La transcription lit « il vient à vérifier que $k$ est
algébriquement fermé dans $A$ », en signalant que les deux mots du milieu sont
serrés et la lecture incertaine. C'est bien une vérification restante et non
une assertion : le feuillet s'arrête là, le reste est vierge.}

\subsection*{La feuille de critères (page 36)}

Un feuillet écrit à l'italienne sur papier faiblement réglé, en deux colonnes
que rien ne relie explicitement, et dont les mots de liaison sont pour la
plupart perdus. Ce sont des \emph{critères de séparabilité} en caractéristique
$p$, et ce qui se lit est ceci. D'un côté, la caractérisation par les
dérivations : pour $A \subset B$,
\[
  \delta f = 0 \iff f \in A[B^{p}] ,
  \qquad
  \delta^{(\infty)} f = 0 \iff f \in A[B^{p^{k}}] ,
\]
c'est-à-dire qu'un élément annulé par toutes les $A$-dérivations de $B$ est
exactement un élément de la sous-algèbre engendrée par les puissances
$p$-ièmes, et qu'un élément annulé par toutes les dérivations
itérées est dans les puissances $p^{k}$-ièmes.\footnote{L'exposant du second
critère est lu $(\infty)$ ; il est écrit très petit et pourrait être un $n$, et
c'est le $B^{p^{k}}$ du second membre qui décide en faveur de l'itération.
L'énoncé demande que $B$ ait une $p$-base sur $A$ — c'est l'hypothèse sous
laquelle la sous-algèbre des constantes des dérivations est exactement
$A[B^{p}]$ — et le feuillet la mentionne quatre lignes plus bas
(« $B/A$ soit une $p$-base »), sans la rattacher explicitement aux deux
équivalences. Le rattachement est nôtre.} De l'autre côté, les critères qui
n'utilisent pas les dérivations : $\Omega^{1}_{B/A} = 0$, la même condition
jointe à la platitude, et une condition sur $\Delta_{B/A}$. Le feuillet attribue
l'un des critères à Chevalley.\footnote{La transcription lit « (Chev. Weyl) »
sans certitude. « Chev. » est vraisemblablement Chevalley ; le second nom n'a
pas été identifié et rien n'en est proposé ici.}

\pagerange{39}{40}
\section*{VII. Foncteurs différentiels d'ordre $n$ (pages 39 et 40)}

Deux feuillets sous la chemise « Foncteurs différentiels d'ordre $n$ », et ce
sont les plus abîmés du premier mouvement : la page 39 est un brouillon rapide
et lourdement biffé, la page 40 est traversée de deux longues diagonales et son
tiers supérieur porte trois couches d'écriture superposées. \emph{Le titre
annonce une définition que le dossier n'atteint pas.} Ce qui suit est donc une
description de ce qui est sur les feuillets, et non un énoncé.

Le cadre posé est celui des catégories fibrées. $S$ est un préschéma,
$\mathcal{E}\ell_{S}$ la catégorie des préschémas sur $S$, $\mathcal{C}$ une
catégorie fibrée sur elle — le feuillet précise « pleine », d'une lecture
douteuse — et $\mathcal{C}'$ une autre. On se donne un foncteur
\[
  F : \mathcal{C} \longrightarrow \Phi^{\mathcal{C}}
\]
dont le feuillet demande qu'il soit une « section continue », et l'on regarde
ses objets $\{X'/S'\}$, $\{X''/S''\}$ au-dessus d'une tour $S \to S' \to S''$,
en exigeant la compatibilité aux extensions de la base. L'exemple donné est
$\underline{P}^{n}_{X'/S'}$, c'est-à-dire précisément le foncteur des parties
principales construit à la section IV.

L'intention se laisse donc deviner — un foncteur entre catégories fibrées est
\emph{différentiel d'ordre $\leqslant n$} lorsqu'il se factorise à travers le
$n$-ième voisinage infinitésimal, ce qui est la forme abstraite de la notion
d'opérateur différentiel d'ordre $\leqslant n$ — mais les feuillets ne
l'écrivent pas, et elle n'est pas écrite ici à leur place. La page 40 ajoute
des schémas auxiliaires $S^{d,n} = \mathrm{Spec}\,\bigl(A[t_{1},\ldots,t_{n}] /
(t_{i} - \cdots)\bigr)$ et $E^{n}_{S^{d,n}} = E^{n} \times_{S} S^{d,n}$, munis
d'une section $\varepsilon : S \to S^{d,n}$, sans que l'anneau quotient se
laisse lire entièrement ni que l'usage qui en est fait apparaisse.

\pagerange{43}{46}
\section*{VIII. Les errata de Lê Dũng Tráng aux §§ 20 et 21 (pages 43 à 46)}

Ici commence le deuxième mouvement, et il n'est plus de sa main. Une lettre
datée « Paris le 11/01/68 », adressée « à Monsieur le Professeur A.
Grothendieck, IHES, Bures s/Yvette », accompagnée de trois feuillets d'errata
numérotés 1, 2 et 3 par leur auteur, chaque chiffre entouré d'un cercle. C'est la pièce que l'inventaire désigne par
« lettres (1968) ».

La lettre elle-même dit exactement ce que les errata sont, et mérite d'être
citée pour ce qu'elle enseigne sur la lecture d'un texte en cours
d'impression : « Les errata ont été relevés pour la plupart dans les
définitions. Je ne sais pas s'ils entravent des démonstrations car celles-ci
faisaient intervenir des éléments de mathématiques que je ne possède pas comme
les idéaux associés ou la profondeur d'un anneau. Pardonnez donc la trivialité
de ces errata. »

Ils ne sont pas triviaux, et le motif qui les traverse est un seul. Le
paragraphe 20 du chapitre IV construit le faisceau des fonctions méromorphes
$\mathcal{K}_X$ en inversant les sections régulières de $\mathcal{O}_X$, et la
difficulté est que \emph{la régularité d'une section n'est pas une propriété
locale au sens naïf} : une section peut être régulière sur un ouvert $U$ et
cesser de l'être par restriction à un ouvert plus petit. Toutes les corrections
de Lê, sauf la dernière, réparent ce point.

\begin{itemize}
  \item \emph{(20.1.3).} $\mathcal{S}(\mathcal{O}_X)$ doit être défini comme le
  sous-faisceau dont les sections sur $U$ sont les éléments réguliers de
  $\Gamma(U, \mathcal{O}_X)$ \emph{dont la restriction à tout ouvert $V \subset
  U$ est encore régulière}. Avec cette clause c'est bien un faisceau, puisqu'un
  élément y est si et seulement si tous ses germes sont réguliers. Il ajoute
  l'avertissement qui est le cœur de l'affaire : $\mathcal{S}(\mathcal{O}_X)_x$
  n'est \emph{pas} nécessairement l'ensemble des éléments réguliers de
  $\mathcal{O}_{X,x}$.
  \item \emph{(20.1.4).} L'anneau $\Gamma(U, \mathcal{K}_X)$ n'est pas l'anneau
  \emph{total} des fractions de $\Gamma(U, \mathcal{O}_X)$, mais un anneau de
  fractions à dénominateurs dans une partie multiplicative d'éléments
  réguliers.
  \item \emph{(20.1.8).} La démonstration de l'inversibilité d'une section de
  $\mathcal{S}(\mathcal{K}_X)$ doit être rédigée localement, avec la même clause
  de stabilité par restriction.
  \item \emph{(20.1.11).} Dans la définition de $\mathcal{S}_f(U)$, lire
  « sections régulières $s \in \Gamma(U, \mathcal{S})$ » et non
  « $s \in \Gamma(U, \mathcal{O}_X)$ ».
  \item \emph{(21.1.5) et (21.1.6).} Les deux mêmes rectifications reportées au
  paragraphe 21.
  \item \emph{(21.2.9).} Deux additions substantielles, qui fournissent la
  démonstration manquante que (21.2.9.1) est un isomorphisme, et qui
  \emph{définissent} la section méromorphe régulière $s_{D}$ de
  $\mathcal{O}_X(D)$ — celle qui correspond à la section $1$ de
  $\mathcal{K}_X$ — pour pouvoir écrire
  \[
    \mathrm{div}(s_{D}) = D .
  \]
  C'est le seul erratum qui ajoute un énoncé plutôt que d'en réparer un.
  \item \emph{(21.2.12).} Ajouter l'hypothèse « localement noethérien » sur le
  préschéma $X$.
  \item \emph{(21.6.2).} Lire $\mathfrak{Z}^{p}(X)$ au lieu de $Z^{p}(X)$ — une
  coquille de fonte.
\end{itemize}

Un dernier erratum était prévu pour (21.3.6) : le numéro est écrit puis barré de
hachures croisées, et rien ne le suit. L'objection a été abandonnée avant
d'être formulée.

\pagerange{48}{48}
\section*{IX. EGA IV 18 ; EGA VI : les sections d'un morphisme non ramifié (page 48)}

Un feuillet isolé, de sa main, à l'encre, rapide, portant ses deux titres — « EGA
IV 18 » en tête, « EGA VI » au milieu, séparés par un trait horizontal. C'est le
feuillet le plus difficile du troisième lot, et plusieurs de ses mots de liaison
sont perdus.

Le second titre mérite qu'on s'y arrête : \emph{le chapitre VI des Éléments n'a
jamais paru}. Ce feuillet est donc un fragment de ce qui devait y aller, et ce
qu'il y met est la descente des sections d'un morphisme non ramifié.

\emph{Premier énoncé (EGA IV 18).} Soit $u : X \to Y$ non ramifié. Soient $U$ un
ouvert schématiquement dense de $Y$, $Y_{0}$ un sous-schéma fermé de $Y$ tel que
$|Y_{0}| = |Y|$, $u_{U}$ une section de $X$ sur $U$, $u_{Y_{0}}$ une section de
$X$ sur $Y_{0}$, coïncidant sur $U_{0} = Y_{0} \cap U$. Alors il existe une
\emph{unique} section $u$ de $X$ sur $Y$ dont les restrictions à $Y_{0}$ et à
$U$ soient $u_{Y_{0}}$ et $u_{U}$.

C'est un énoncé de recollement : un ouvert schématiquement dense et un
sous-schéma fermé de même espace sous-jacent \emph{couvrent} $Y$ pour ce qui
concerne les sections d'un morphisme non ramifié. La non-ramification est ce qui
donne l'unicité — deux sections qui coïncident sur un sous-schéma
schématiquement dense coïncident.

\emph{Second énoncé (EGA VI).} Soit $X \to Y$ non ramifié, $Y' \to Y$ localement
noethérien, fini et épimorphisme, et $Y'' \to Y' \times_{Y} Y'$ surjectif. Alors
la suite
\[
  X(Y) \longrightarrow X(Y') \rightrightarrows X(Y'')
\]
est exacte : le foncteur des points de $X$ est un faisceau pour cette
topologie.\footnote{Le mot « noeth. » sous la flèche $Y' \to Y$ est surchargé et
se lit comme biffé — ce qui serait cohérent avec la note marginale qui suit —
mais la biffure n'est pas certaine et la transcription le laisse dans le
diagramme.} La marge de droite porte, en vertical, le dernier mot souligné :
« On doit pouvoir éliminer l'hyp. loc. noeth. et remplacer fini par
\emph{entier} ». C'est le programme que le chapitre VI aurait exécuté.

Suit un \emph{corollaire des deux} — la variante du premier énoncé où $Y_{0}$
est remplacé par un $Y'$ fini surjectif sur $Y$ — et un second corollaire, « Si
$f : X \to Y$ est … », que le feuillet interrompt sur ses propres points de
suspension.\footnote{Les deux mots entre « est » et les points de suspension
n'ont pas pu être lus ; la transcription note que le premier pourrait être
« un. » et que le second commence par un f. Aucune reconstitution n'est proposée
ici : l'énoncé du corollaire est perdu.}

\pagerange{49}{53}
\section*{X. Suites régulières : la définition de 16.9.2 n'est pas la bonne (pages 49 à 53)}

C'est le sommet du dossier, et la seule de ses questions qui se pose, se
travaille et se tranche à l'intérieur de la même chemise. Elle porte, de sa
main, en haut à droite d'un feuillet par ailleurs blanc : « Suites régulières /
Th. 11.2.9 de Raynaud ».

\subsection*{Les deux fiches de travail (pages 50 et 51)}

Deux feuillets à en-tête de l'IHES, et ce sont des listes. Le premier, au crayon,
inventorie les \emph{applications} du théorème 11.2.9 dans la suite du chapitre
IV : 11.3.4 (i), 11.3.8 b'), 17.12.1 c'), 19.2.4 c'), 19.4.6 (ii), et une note
soulignée, « À signaler » suivi d'un mot court qui ne se lit pas, puis
« remarques ? ». Le second, à l'encre, inventorie
les paragraphes où la notion est en jeu : $0_{\mathrm{IV}}$ 15, IV 11.3.7 et
11.3.8, IV 16.9, IV 19, IV 21.15 — puis, sans transition, la ligne qui fait de
cette fiche autre chose qu'un index :

\begin{quote}
Lettres Deligne 17. Oct. 66. et 29 Oct. 66.
\end{quote}

et, au-dessous, encadré de deux traits horizontaux et barré de deux longues
diagonales :

\begin{quote}
Notion d'immersion régulière dans le cas non noethérien : est-elle
vraisemblable ?? Vérifier § 16 !!
\end{quote}

Les lettres qu'il cite sont dans le dossier, quatre feuillets plus loin, et la
mathématique qui répond à la question est dedans. La question est rayée parce
qu'elle a reçu sa réponse, et cette réponse est le tapuscrit qui suit
immédiatement.

\subsection*{Le post-scriptum : la notion n'est pas la bonne (pages 52 et 53)}

Deux feuillets de tapuscrit, copie carbone, sans en-tête ni signature, à la
première personne et de sa voix — « Dans mon séminaire sur Riemann-Roch cette
année ». Le destinataire n'est pas nommé. C'est un post-scriptum de lettre, et
il est complet sur les deux pages.

L'aveu tient en une phrase : « il apparait (hélas) que la notion introduite sous
ce nom dans EGA 16.9.2. n'est pas la bonne notion ». Et la raison tient en une
autre, qui est le genre d'observation dont dépend la solidité d'un chapitre
entier :

\emph{Si $A$ est un anneau local et $x$, $y$ deux éléments de son idéal maximal,
il est facile de donner un exemple où $(x,y)$ est une suite régulière et
$(y,x)$ ne l'est pas.}

Une notion qui dépend de l'ordre dans lequel on écrit les générateurs ne jouit
d'aucune des propriétés qu'on en attend : on ne peut pas prendre une base de
$J/J^{2}$ et la relever, la notion d'immersion régulière ne se comporte pas bien
par localisation plate ni même étale, et « toute sorte d'autres » inconvénients
s'ensuivent.

Ce qu'il propose à la place est la condition de \emph{Koszul}. Une suite
$(f_{1},\ldots,f_{n})$ d'éléments d'un anneau, engendrant un idéal $J$, est
\emph{Koszul-régulière} lorsque le complexe de l'algèbre extérieure construit
sur elle est acyclique en degrés $> 0$.\footnote{Le tapuscrit écrit « que le
complexe de l'algèbre extérieure construit avec soit nul en degrés $0$ », le
nom du complexe manquant après « avec » et l'inégalité étant frappée sans signe.
La condition voulue est l'acyclicité en degrés strictement positifs, c'est-à-dire
que le complexe de Koszul $K_{\bullet}(f, A)$ soit une résolution de $A/J$ ;
c'est ainsi qu'elle est écrite ici, et c'est la condition (iv) que Deligne
démontre équivalente aux autres à la page 62.} Il énonce alors ce qui rend la
notion bonne, et c'est exactement ce qui manquait à l'autre :

\begin{itemize}
  \item une suite Koszul-régulière est \emph{quasi-régulière} ;
  \item \emph{toute autre suite de $n$ générateurs de $J$} a la même propriété —
  la condition ne dépend donc pas du système de générateurs choisi, ni de leur
  ordre ;
  \item les trois notions qui en découlent — suite K-régulière, idéal régulier,
  immersion régulière — \emph{se descendent bien par fpqc} ;
  \item un composé de deux immersions régulières en est une, et les deux autres
  propriétés de transitivité de 19.1.5 (iii) « semblent marcher encore sans
  hypothèse noethérienne », Berthelot en ayant déjà vérifié une ;
  \item et pour ce qu'il veut en faire dans le séminaire Riemann-Roch, « c'est
  bien la nouvelle notion d'immersion régulière qui fait tout marcher ».
\end{itemize}

\subsection*{Ce qu'il décide d'en faire}

La seconde page est la décision pratique, et c'est elle qui fait de ce
tapuscrit un document et non une note. Il est trop tard pour corriger les
épreuves de IV 16 et 19. On ajoutera donc une remarque 19.1.8 « dans laquelle
nous ferions l'aveu » que la définition adoptée s'est avérée peu maniable dans
le cas non noethérien, en renvoyant le lecteur, pour les détails et en attendant
la réédition du chapitre IV, au séminaire SGA 6 « qui devrait sortir vers la fin
de l'année 67 ou début 68 » ; et l'on insérera, au paragraphe 16, une note de
bas de page renvoyant à cette remarque. Il ajoute entre parenthèses ce qui
limite les dégâts : « Dans le par. 16, c'est en fait la notion quasi-régulière
qui sert vraiment ».

Puis vient la proposition de terminologie, et c'est la page la plus instructive
du dossier sur la manière dont un vocabulaire se fixe. À la réédition il faudra
revoir $0_{\mathrm{IV}}$ 15 dans le même esprit, et y remplacer partout
l'hypothèse « suite régulière » par l'hypothèse plus faible et plus maniable
« suite Koszul-régulière ». Et alors :

\begin{quote}
il s'imposera alors, pour les suites également, de réserver le mot le plus court
à la notion qui manifestement est la plus serviable, et il faudra envisager de
changer notre terminologie, en disant « suite régulière » au lieu de « suite
Koszul-régulière », et en réservant tout au plus un terme tel que « suite
strictement régulière » à la notion éculée.
\end{quote}

Il termine par la remarque qui explique pourquoi rien de tout cela n'avait été
vu plus tôt : « Bien sûr, dans le cas noethérien, les trois variantes
(régulière, quasi, et Koko) sont équivalentes. »\footnote{« Koko » pour
« Koszul » est de la frappe. Le tapuscrit porte aussi, frappé, « il faudra
envisaher » pour « envisager », et une marge gauche au crayon souligne d'un
double trait vertical tout le passage allant de « ajouter une remarque 19.1.8 »
à « la définition incriminée » : la marque est de lui, et c'est l'endroit qu'il
retient.}

La décision a été suivie pour moitié. La condition de Koszul est bien celle
dont on se sert hors du cas noethérien, et c'est elle qui gouverne la notion
d'immersion régulière telle qu'on l'emploie aujourd'hui. Mais l'échange de noms
qu'il envisage n'a pas eu lieu : « suite régulière » a gardé son sens ancien, et
c'est la notion serviable qui porte le qualificatif, sous le nom de \emph{suite
Koszul-régulière}. La réédition du chapitre IV, dont tout ce feuillet suppose
la venue, n'a pas paru.

\pagerange{55}{64}
\section*{XI. Les lettres de Deligne (pages 55 à 64)}

Une chemise par ailleurs blanche, la page 54, porte de sa main les deux mots
« Lettres Deligne ». Ce qu'elle couvre est dix feuillets à l'encre, de la main
de Pierre Deligne, annotés au crayon par le destinataire.

Deux avertissements sont nécessaires avant de les lire, et ni l'un ni l'autre ne
peut se voir à l'intérieur d'un seul lot de transcription.

\emph{Il y a trois lettres et non quatre.} Les feuillets 61 à 64, qui suivent
immédiatement, sont paginés 3 à 6 en haut à droite par leur auteur et signés
« P. Deligne » : ce sont les pages 3 à 6 de la lettre du 17 octobre 1966, dont
les pages 1 et 2 sont les feuillets 59 et 60. La soudure est exacte — la page 60
s'achève en construisant $A' = A[[T_{i}]]$ et en établissant l'isomorphisme
$\mathfrak{m}'/(\mathfrak{m}'^{2} + \mathfrak{m}A') \simeq
\mathfrak{n}/(\mathfrak{n}^{2} + \mathfrak{m}B)$ ; la page 61 s'ouvre sur
« Remplaçant $A$ par $A'$, il reste à prouver que si
$\mathfrak{m}/\mathfrak{m}^{2} \otimes K \to \mathfrak{n}/\mathfrak{n}^{2}$
est surjectif… », qui est exactement le point où l'on vient d'arriver.\footnote{La
soudure se vérifie une seconde fois, et indépendamment, par la fin de la
lettre : à la page 64, Deligne écrit « je suis honteux de la démonstration
tordue donnée pg 2 de ce que $\mathrm{Tor}$ et “pro-objet complétion”
commutent ». Sa page 2 est le feuillet 60 du dossier, et la proposition sur la
commutation de $\mathrm{Tor}$ avec la complétion pro-objective y est en effet.
Les deux transcriptions disent chacune que la lettre se poursuit dans l'autre
lot ; aucune ne pouvait établir la jonction.}

\emph{Elles sont classées dans l'ordre inverse des dates.} Le dossier donne 15
décembre, puis 29 octobre, puis 17 octobre. Lues ainsi, elles sont
incompréhensibles : la lettre du 15 décembre répond à une lettre du 28 et
reprend une objection formulée le 29 octobre. Elles sont donc lues ici dans
l'ordre où elles ont été écrites, ce qui oblige à remonter les pages — d'abord
59 à 64, puis 57 et 58, puis 55 et 56.

Les annotations au crayon dans les marges sont de Grothendieck. Elles sont
courtes — « oui », « Err », « ! », « ? », « Th. de Auslander », « dans réédition
du Chap I ? » — et elles sont son seul mot dans cet échange : c'est le dossier
d'un homme qui reçoit.

\pagerange{59}{64}
\subsection*{a) La lettre du 17 octobre 1966 (pages 59 à 64)}

Six pages, la plus longue et la plus riche. Elle s'ouvre sur deux phrases qu'il
faut garder telles quelles : « Je n'ai pas l'intention d'essayer de démontrer,
cette année, le “théorème de Lefschetz vache”. J'espère donc que tu trouves
quelqu'un qui le fasse. » Et : « J'ai été malade au début du mois et, comme
lecture de convalescence, je lis EGA IV. »

\emph{$0_{\mathrm{IV}}$ 19.3.12 est faux.}

Le premier verdict est net. Il faut prendre pour $W_{\lambda}$ l'ensemble des
homomorphismes d'\emph{algèbre} $B \to C_{\lambda}$ relevant $B \to C_{\lambda}
/ J_{\lambda}$, et \emph{ce n'est pas une variété linéaire}, puisqu'on n'a pas
$J_{\lambda}^{2} = 0$. Pour sauver l'argument il faudrait savoir que, pour $A$
artinien, les parties de $A^{I}$ définies par des équations $P(x_{i}) \in
\mathfrak{b}$ jouissent d'une propriété de compacité — ce qui est faux dès que
$I$ est infini, et déjà pour $A$ un corps infini : prenant $I = k \cup
\{\infty\}$ et les parties $(x_{\infty} - a)\,y_{a} = 1$, les intersections
finies sont non vides et l'intersection totale est vide. Il conclut : « Je crois
qu'il n'y a pas d'espoir. » La marge porte, au crayon et entouré, un seul mot :
\emph{oui}.

\emph{La formule d'Auslander (page 59).}

À propos de $0_{\mathrm{IV}}$ 17.3.4, il signale « une jolie démonstration
cohomologique de ce résultat, qui fournit un résultat plus complet (lequel se
trouve dans Northcott) », et la marge au crayon le nomme : \emph{Th. de
Auslander}.

\emph{Proposition.} Soient $A$ un anneau local noethérien et $M$ un $A$-module
de type fini, non nul, de dimension projective finie. Alors
\[
  \dim\mathrm{proj}(M) = \mathrm{prof}\,A - \mathrm{prof}\,M .
\]

C'est la formule d'Auslander-Buchsbaum, et la démonstration de la lettre est
celle qu'on donne aujourd'hui. On part de la formule, valable pour un complexe
parfait $L$,
\[
  R\mathrm{Hom}(K, L) = R\mathrm{Hom}(K, A) \overset{L}{\otimes} L ,
\]
avec $K = k$ et $L = M$ ; on en tire la suite spectrale
\[
  E_{2}^{pq} = \mathrm{Tor}_{-p}\bigl(\mathrm{Ext}^{q}(k, A),\, M\bigr)
  \Longrightarrow \mathrm{Ext}^{p+q}(k, M) .
\]
Posant $d = \dim\mathrm{proj}\,M$ et $a = \mathrm{prof}\,A$, le terme $E_{2}$ est
nul pour $q < a$ et pour $p < -d$ : c'est une bande, et Deligne la dessine en
marge.\footnote{Le « dessin ci-contre » est, en bas à gauche du feuillet : un
axe horizontal $p$ et un axe vertical $q$, une bande verticale hachurée de
largeur $d$ mesurée à gauche de l'axe des $q$, commençant à la hauteur $a$
au-dessus de l'axe des $p$. C'est le support de $E_{2}^{pq}$, et c'est le
dessin qui contient la démonstration : le coin de la bande est le seul terme
qui survive.} Le coin de la bande donne alors
\[
  \mathrm{Ext}^{n}(k, M) = 0 \text{ pour } n < a - d ,
  \qquad
  \mathrm{Ext}^{a-d}(k, M) = \mathrm{Tor}_{d}\bigl(\mathrm{Ext}^{a}(k, A),\,
  M\bigr) \neq 0 ,
\]
la non-nullité venant de ce que $\mathrm{Ext}^{a}(k,A)$ est somme de copies de
$k$ et que $\mathrm{Tor}_{d}(k, M) \neq 0$ par définition de $d$. Donc
$\mathrm{prof}\,M = a - d$.

La page 60 étend le résultat : il vaut encore, en remplaçant $\dim\mathrm{proj}$
par la $\mathrm{Tor}$-dimension, pour $M$ de type fini sur un anneau noethérien
$B$ qui est une $A$-algèbre avec $\mathfrak{m}B \subset \mathrm{rad}(B)$.

\emph{$\mathrm{Tor}$ et la complétion pro-objective (page 60).}

L'extension repose sur un énoncé pour lequel, dit-il, « je ne connais pas de
référence ».

\emph{Proposition.} Soient $A$ un anneau local noethérien, $B$ une $A$-algèbre
locale noethérienne ($A \to B$ local), $\mathfrak{m}$ un idéal de $A$, $N$ un
$A$-module de type fini et $M$ un $B$-module de type fini. Alors
\[
  \varprojlim_{k} \mathrm{Tor}_{p}^{A}(N/\mathfrak{m}^{k}N,\, M)
  = \mathrm{Tor}_{p}^{A}(N, M)^{\wedge} ,
\]
et plus précisément on a une égalité de \emph{pro-objets}
\[
  \varprojlim\;\mathrm{Tor}_{p}^{A}(N, M) \big/ \mathfrak{m}^{k}\,\mathrm{Tor}_{p}^{A}(N, M)
  \;\xrightarrow{\;\sim\;}\;
  \varprojlim\;\mathrm{Tor}_{p}^{A}(N/\mathfrak{m}^{k}N,\, M) .
\]

Si $B$ était plat sur $A$, on prendrait une résolution de $M$ sur $B$ et
Artin-Rees conclurait — c'est lui qui exprime que « le pro-objet complétion »
est un foncteur exact. On se ramène donc, en remplaçant $B$ par une algèbre dont
il est quotient et en le complétant, à un lemme de structure :

\emph{Proposition.} Soit $A \to B$ un homomorphisme local d'anneaux locaux
noethériens, avec $B$ complet. Alors $B$ est quotient d'un anneau $C$ tel que
(i) $C$ soit local et $A \to C$, $C \to B$ locaux ; (ii) $C$ soit \emph{plat sur
$A$} ; (iii) si $k$ est le corps résiduel de $A$, $C_{0} = C \otimes_{A} k$ soit
\emph{régulier} et $\mathrm{gr}\,C_{0} \to \mathrm{gr}\,B_{0}$ bijectif en
degrés $0$ et $1$.\footnote{Dans (iii), le mot « régulier » est écrit par Deligne
au-dessus d'un passage biffé dont on lit encore « de séries formelles sur une
extension de $k$ » : il a affaibli sa propre condition, et c'est la version
corrigée qui est retenue. Le « et et » de la condition (i) est ainsi sur le
feuillet. La marge porte, au crayon et entouré, en oblique, la question de
Grothendieck : « dans réédition du Chap I ? ».}

\emph{Le relèvement d'une extension résiduelle purement inséparable (pages 60 à 62).}

Ici commence la démonstration, et c'est elle qui franchit la limite des lots. On
choisit une base $t_{i}$ d'un supplémentaire de l'image de
$\mathfrak{m}/\mathfrak{m}^{2} \otimes_{k} K \to
\mathfrak{n}/\mathfrak{n}^{2}$, on la relève dans $B$, ce qui définit
$A[[T_{i}]] \to B$, et l'on pose $A' = A[[T_{i}]]$, plat sur $A$, avec
$\mathfrak{m}'/(\mathfrak{m}'^{2} + \mathfrak{m}A') \xrightarrow{\sim}
\mathfrak{n}/(\mathfrak{n}^{2} + \mathfrak{m}B)$. Remplaçant $A$ par $A'$, il
reste à traiter le cas où $\mathfrak{m}/\mathfrak{m}^{2} \otimes K \to
\mathfrak{n}/\mathfrak{n}^{2}$ est surjectif, c'est-à-dire $K = B \otimes_A k$,
et il s'agit alors de trouver $C$ avec $K = C \otimes_A k$. « On va manger $K$
par étapes », écrit-il, et les trois étapes sont :

\begin{enumerate}
  \item[a)] \emph{La partie transcendante.} Soit $(x_{i})_{i \in I}$ une base de
  transcendance de $K$ sur $k$ ; on relève les $x_{i}$ dans $B$, définissant
  $A(x_{i}) \to B$, où $A(x_{i})$ est la limite inductive des localisés, au
  point générique de la fibre spéciale, des $A[X_{i}]_{i \in J}$ pour $J$ fini
  dans $I$ — la construction de Nagata.\footnote{La lettre renvoie à Nagata,
  \emph{Local rings}, p. 17 et 18, et indique deux manières de conclure à la
  noethérianité ou à la platitude : $0_{\mathrm{III}}$ 10.3.1.3, ou bien
  compléter et appliquer le critère de platitude.}
  \item[b)] \emph{La partie séparable.} On est ramené à $K$ algébrique sur $k$,
  et l'on trouve « sans histoire » $A'$ plat sur $A$, non ramifié, de corps
  résiduel la clôture séparable de $k$ dans $K$.
  \item[c)] \emph{La partie purement inséparable}, « et on va utiliser une
  propriété “d'intersection complète” des corps, que je n'ai pas su bien
  dégager ».
\end{enumerate}

Le point c) est le cœur. Soit $K_{n} = K \cap k^{p^{-n}}$ le sous-corps des $x$
tels que $x^{p^{n}} \in k$, et $x_{n}^{\alpha}$ une $p$-base de $K_{n}$ sur
$K_{n-1}$. Alors $K$ est engendré par les $x_{n}^{\alpha}$ soumis aux relations
$(x_{n}^{\alpha})^{p} = P_{n}^{\alpha}(x_{m}^{\beta},\ m < n)$ à coefficients
dans $k$. Relevant les $x_{n}^{\alpha}$ dans $B$ et corrigeant par récurrence,
on obtient
\[
  \boxed{\;(x_{n}^{\alpha})^{p} = P_{n}^{\alpha}(x_{m}^{\beta},\ m < n)
  + \sum_{i \geqslant 1} Q_{n,i}^{\alpha}(x_{m}^{\beta}) \;}
\]
où $P_{n}^{\alpha}$ est à coefficients dans $A$ et $Q_{n,i}^{\alpha}$ à
coefficients dans $\mathfrak{m}^{i}$.\footnote{L'exposant de $\mathfrak{m}$ est
entouré d'un cercle sur le feuillet et lu $(n)$ par le transcripteur ; en marge
de gauche, au crayon et d'une autre main que celle de la lettre — celle du
destinataire — on lit « $Q_{n i}^{\alpha}$ linéaire en les $x_{m}^{\beta}$ » et
« $(\ )^{i}$ ? ». L'annotateur a raison et la transcription le note : l'exposant
attendu est $i$, non $n$, puisque la série $\sum_{i} Q_{n,i}^{\alpha}$ doit
converger $\mathfrak{m}$-adiquement. C'est le seul endroit du dossier où une
marginale de Grothendieck corrige un calcul de Deligne, et la correction est
juste.} On prend alors pour $C$ l'anneau complet engendré par des générateurs
$x_{n}^{\alpha}$ soumis à ces relations — rigoureusement, la limite projective
des $A/\mathfrak{m}^{k}$-algèbres définies par les relations tronquées à $i
\leqslant k-1$.

Reste à voir que c'est plat, et la vérification est courte : il suffit de le
faire pour le quotient de $A[X_{n}^{\alpha}]$ par un nombre fini de relations,
où seules un nombre fini d'indéterminées interviennent ; on applique
$0_{\mathrm{IV}}$ 15.1.16 c) $\Rightarrow$ b), et il faut vérifier que dans
$k[X_{n}^{\alpha}]$ les $(X_{n}^{\alpha})^{p} -
P_{n}^{\alpha}(X_{m}^{\beta},\ m<n)$ forment une suite régulière, « et cela
résulte d'un calcul de dimension trivial — CQFD ».

Il ajoute la formulation générale qu'il n'a pas su dégager, et qui vaut d'être
retenue : \emph{toute extension algébrique $K/k$ est du type $K =
k[T_{\alpha}]/(P_{i})$ de sorte qu'en relevant les $P_{i}$ à un anneau artinien
$A$ de corps résiduel $k$, l'algèbre $A[T_{\alpha}]/(\widetilde{P}_{i})$ soit
plate sur $A$.} C'est une propriété « d'intersection complète » des corps, et
c'est ce qui fait marcher tout le point c).

\emph{Les résultats manquants sur les suites régulières (pages 62 et 63).}

« Il manque des résultats sur les suites régulières, notamment le suivant auquel
Quillen faisait une référence fantôme dans sa lettre sur $\Omega_{B/A}$. » La
marge du destinataire, en oblique et d'une encre plus grasse, répond :
« garder EGA IV 19, à compléter éventuellement. » Et une autre, au-dessus :
« $0_{\mathrm{IV}}$ 15 », le $16.4$ qui la précédait étant biffé.

C'est ici que le dossier répond à la question rayée de la page 51.

\emph{Proposition.} Soient $A$ noethérien et $I$ un idéal. Les conditions
suivantes sont équivalentes :

\begin{enumerate}
  \item[(i)] $\mathrm{Sym}_{A/I}(I/I^{2}) \xrightarrow{\sim} \mathrm{gr}_{I}(A)$
  et $I/I^{2}$ est projectif sur $A/I$ ;
  \item[(ii)] $\Lambda^{\bullet}(I/I^{2}) \xrightarrow{\sim}
  \mathrm{Tor}^{A}_{\bullet}(A/I, A/I)$ et \emph{idem}.
\end{enumerate}

Si de plus $A$ est local et si $f_{1},\ldots,f_{n}$ est un système
\emph{minimal} de générateurs de $I$ — c'est-à-dire $n = \dim_{k} I \otimes k$ —
cela équivaut encore à

\begin{enumerate}
  \item[(iii)] $f_{1},\ldots,f_{n}$ est une suite régulière ;
  \item[(iv)] le complexe de l'algèbre extérieure $K(f, A)$ est une résolution
  de $A/I$ ;
  \item[(v)] $\mathrm{Ext}^{i}(A/I, A) = 0$ pour $i < n$.
\end{enumerate}

L'hypothèse de \emph{minimalité} du système de générateurs est ce qui rend
l'équivalence vraie et ce qui en marque la limite : c'est précisément le passage
à un système quelconque de générateurs qui casse dans le cas non noethérien, et
c'est ce que la notion de Koszul répare.

La démonstration donnée est celle-ci. On se ramène à $A$ local et $I \neq A$.
Les implications (i) $\Leftrightarrow$ (iii) $\Rightarrow$ (iv) $\Rightarrow$
(ii) et (iv) $\Rightarrow$ (v) sont claires, la dernière parce que (iv) entraîne
que $\mathrm{Tor}^{1}(A/I, A/I) = I/I^{2}$ est $A/I$-libre. Pour (ii)
$\Rightarrow$ (iv) : soit $P_{*}$ une résolution projective de $A/I$ ; il existe
une seule classe d'homotopie d'applications $K(f,A) \to P_{*}$ induisant
l'identité sur $A/I$, et l'application déduite $\Lambda^{\bullet}(I/I^{2}) =
H_{*}(K(f,A) \otimes A/I) \to \mathrm{Tor}(A/I,A/I) = H_{*}(P_{*} \otimes A/I)$
coïncide avec celle de (ii). Soit $E$ le cylindre de cette application : (ii)
signifie que $E \otimes A/I$ est acyclique, donc par Nakayama que $E$ l'est,
c'est-à-dire (iv).

Reste (v) $\Rightarrow$ (iii), et Deligne l'ajourne d'une ligne — « Va être
modulée » — puis démontre une proposition plus forte, à coefficients dans un
module :

\emph{Proposition.} Si $I$ est un idéal de l'anneau local $A$, $f_{1},\ldots,
f_{n}$ un système minimal de générateurs de $I$, $M$ un $A$-module de type fini,
et si $\mathrm{Ext}^{i}(A/I, M) = 0$ pour $i < n$, alors $f_{1},\ldots,f_{n}$
est une suite $M$-régulière.

La démonstration est en trois temps, dont le troisième est le joli. On remarque
d'abord que la conclusion ne dépend qu'apparemment du système de générateurs, en
déroulant la chaîne d'équivalences entre régularité, quasi-régularité et
bijectivité de $\mathrm{gr}_{I}(A) \otimes_{A/I} M/IM \to \mathrm{gr}_{I}(M)$.
Puis, s'il existe $f \in I \setminus \mathfrak{m}I$ qui soit $M$-régulier, on
descend d'un cran et l'on conclut par récurrence sur $n$. Enfin, s'il n'en
existe aucun, alors
\[
  I \subset \mathfrak{m}I \cup \bigcup_{\mathfrak{p} \in \mathrm{Ass}(M)}
  \mathfrak{p} ,
\]
et un lemme d'évitement conclut qu'il existe $\mathfrak{p} \in \mathrm{Ass}\,M$
avec $I \subset \mathfrak{p}$, d'où $\mathrm{Hom}(A/I, M) \neq 0$, donc $n = 0$,
$I = 0$, et l'assertion est vraie.

\emph{Lemme} (que Deligne attribue par un « (Nagata !) »). Soient $\mathfrak{a}$
un idéal d'un anneau $A$, $(\mathfrak{b}_{i})_{0 < i \leqslant n}$ des idéaux
et $(\mathfrak{p}_{i})_{0 < i \leqslant m}$ des idéaux premiers. Si
$\mathfrak{a} \subset \bigcup_{i} \mathfrak{b}_{i} \cup \bigcup_{i}
\mathfrak{p}_{i}$, alors ou bien $\mathfrak{a} \subset \bigcup
\mathfrak{b}_{i}$, ou bien il existe $i$ avec $\mathfrak{a} \subset
\mathfrak{p}_{i}$.

C'est le lemme d'évitement des idéaux premiers, dans la forme forte où les
idéaux non premiers sont traités en bloc. Sa démonstration en deux lignes est
celle de la lettre : on se ramène à $m = 1$ et à $\mathfrak{p} \not\supset
\mathfrak{b}_{i}$ ; si $\mathfrak{a} \not\subset \bigcup \mathfrak{b}_{i}$, on
prend $x \in \mathfrak{a} \setminus \bigcup \mathfrak{b}_{i}$, donc $x \in
\mathfrak{p}$, et pour $y \in \mathfrak{a} \cap \bigcap \mathfrak{b}_{i}$ on a
$x + y \in \mathfrak{a} \setminus \bigcup \mathfrak{b}_{i}$, donc $x+y \in
\mathfrak{p}$ et $y \in \mathfrak{p}$ ; d'où $\mathfrak{a} \cap \bigcap
\mathfrak{b}_{i} \subset \mathfrak{p}$ et, comme $\mathfrak{b}_{i} \not\subset
\mathfrak{p}$, $\mathfrak{a} \subset \mathfrak{p}$.\footnote{Suit, entièrement
biffé, un N.B. où Deligne croit qu'un dévissage analogue permettrait de montrer
(iv) $\Rightarrow$ (iii) sans l'hypothèse de minimalité. La biffure est la
sienne, et l'énoncé n'est pas démontré dans le dossier. C'est exactement le
point que le tapuscrit des pages 52 et 53 tranchera autrement, en changeant de
définition plutôt qu'en affaiblissant l'hypothèse.}

\emph{Plat $\Longleftrightarrow$ formellement projectif (pages 63 et 64).}

En marge, « $0_{\mathrm{IV}}$ 19.2 », et : « Il manque la proposition suivante,
dont tu démontres un cas particulier en 19.7.1.1 ».

\emph{Proposition.} Soient $A \to B$ un homomorphisme local d'anneaux locaux
noethériens et $M$ un $B$-module de type fini, $A$ et $M$ munis des topologies
$\mathfrak{m}$- et $\mathfrak{n}$-adiques. Alors $M$ est \emph{plat} sur $A$ si
et seulement s'il est \emph{formellement projectif} sur $A$.

La démonstration est un enchaînement d'équivalences dont l'intérêt est la
méthode — tout se passe dans les catégories de pro- et d'ind-objets, et la
dualité entre les deux fait le travail. On se ramène à $A$ artinien local, puis
\begin{align*}
  \text{(ii)} &\iff \varinjlim_{k} \mathrm{Ext}^{1}_{A}(M/\mathfrak{n}^{k}M, N)
  = 0 \quad \forall N \\
  &\iff \text{idem pour } N \text{ un } k\text{-espace vectoriel (dévissage)} \\
  &\iff \varinjlim \mathrm{Ext}^{1}_{A}(M/\mathfrak{n}^{k}M, k) = 0 \text{ comme Ind-objet} \\
  &\iff \varprojlim \mathrm{Tor}^{A}_{1}(M/\mathfrak{n}^{k}M, k) = 0
  \text{ comme pro-objet,}
\end{align*}
le précédent étant le dual de celui-ci. On calcule ce $\mathrm{Tor}$ par une
résolution plate de $k$, et Artin-Rees identifie le pro-objet à
$\varprojlim\ \mathrm{Tor}^{A}_{1}(M,k)/\mathfrak{n}^{k}
\mathrm{Tor}^{A}_{1}(M,k)$, dont la nullité signifie $\mathrm{Tor}_{1}(M,k) = 0$,
c'est-à-dire la platitude.\footnote{L'indice des $\mathrm{Ext}$ et $\mathrm{Tor}$
est un simple trait sur la feuille ; il est lu $1$, et c'est le seul qui donne
un énoncé vrai.}

\emph{Les six dernières remarques, et la signature (page 64).}

Le reste de la lettre est une suite de notes marginales numérotées par le
paragraphe qu'elles visent, et elles montrent le genre de lecture dont un texte
en épreuves a besoin.

\begin{itemize}
  \item \emph{$0_{\mathrm{IV}}$ 16.1.9} : on peut y supprimer l'hypothèse
  noethérienne, en référant à Bourbaki, \emph{Algèbre commutative}, ch. V, § 1,
  n\textsuperscript{o} 1, th. 1.
  \item \emph{$0_{\mathrm{IV}}$ 19.4.6} : il n'est pas nécessaire de supposer
  $B$ formellement projectif dans l'hypothèse 2\textsuperscript{o} ; et dans
  1\textsuperscript{o}, cette hypothèse peut se remplacer par une hypothèse de
  finitude.
  \item \emph{$0_{\mathrm{IV}}$ 19.4.7} : la démonstration se simplifierait si
  l'on avait fait remarquer auparavant que $\mathrm{Exalcotop}_{A}(B, L)$ est
  foncteur semi-exact en $L$.
  \item \emph{$0_{\mathrm{IV}}$ 19.8.3} : lire « homomorphisme local » et non
  « homomorphisme », aux trois endroits indiqués.
  \item \emph{$0_{\mathrm{IV}}$ 20.4.6} : \emph{le signe de $dx$ est le
  contraire de celui que choisit l'analyse}, où $\Delta f(x;x') = f(x') - f(x)$,
  compte tenu de ce que $B \otimes B$ et $P^{1}_{B/A}$ sont des $B$-algèbres par
  $b \mapsto b \otimes 1$. La marge du destinataire, en oblique : « faire
  errata : remplacer $P_{1} - P_{2}$ par $P_{2} - P_{1}$ ! attention aux
  passages subséquents ».
  \item \emph{$0_{\mathrm{IV}}$ 20.4.13 (iv)} : l'hypothèse « $A$ intégralement
  clos » n'a rien à faire ici.
  \item \emph{III 4.3.6} : « la formule $K' = K \otimes_{A} \widehat{A}$ est
  fausse en général ! » — et la marge porte trois points d'exclamation.
\end{itemize}

Et la lettre se ferme sur un aveu : « En relisant cette lettre, je suis honteux
de la démonstration tordue donnée pg 2 de ce que $\mathrm{Tor}$ et “pro-objet
complétion” commutent. Il suffirait, par dévissage, de traiter directement le
cas $N = A$ et $M = B$, ié “$\varprojlim$”
$\mathrm{Tor}^{A}_{p}(A/\mathfrak{m}^{k}, M) = 0$ si $p > 0$, mais j'en suis
incapable. Tu trouveras ci-joint une lettre de Serre que je te renvoie. »
La lettre de Serre n'est pas dans le dossier.

\pagerange{57}{58}
\subsection*{b) La lettre du 29 octobre 1966 (pages 57 et 58)}

Deux pages, sur EGA IV 19. Elles portent sur une seule chose : \emph{les
hypothèses de séparation}.

\emph{L'hypothèse $(*)$ « semble idiote »}

L'hypothèse mise en cause est : « tout quotient de tout sous-module de $M$
(resp. $M^{n}$) est séparé pour la topologie $J$-adique », et l'objection est
en deux points.

\begin{enumerate}
  \item[(i)] Elle n'est \emph{pas impliquée} par l'hypothèse de finitude
  relative noethérienne standard — « $M$ fini sur la $A$-algèbre noethérienne
  $B$, $A$ noethérien, $JB \subset R(B)$ » — « qui elle est naturelle ». La
  marge porte un point d'interrogation au crayon.
  \item[(ii)] Le seul exemple non trivial indiqué où l'hypothèse soit vérifiée,
  soit 19.5.6 (ii), est \emph{faux} : dans $(*)$ on ne se limite pas à prendre
  des quotients par des sous-modules \emph{gradués}. Contre-exemple : $A =
  B[T]$, $M = N[T]$, $N$ un $A$-module où $T$ agit par $1$ ; ce $A$-module est
  quotient de $M$. La marge porte « Err », souligné, au crayon — verdict
  accepté — et un point d'exclamation contre l'exemple.
\end{enumerate}

Suit la phrase qui dit le parti pris : « Je crois qu'il faut éviter d'employer
des hypothèses de séparation qu'on ne rencontre pas dans la nature, comme
“idéalement séparé” etc., et se limiter à utiliser (i) — J'avoue toutefois, avec
honte, que dans le même geste sacrilège je serais prêt à supposer tous les
schémas séparés et noethériens… » Il propose tout de même, pour coiffer (i) et
$(*)$, une hypothèse intermédiaire :
\[
  \forall f : N \to P ,\ N \text{ et } P \text{ finis sur } A ,\quad
  \mathrm{Ker}(M \otimes N \to M \otimes P) \text{ est } J\text{-séparé} ,
\]
contre laquelle la marge porte un point d'exclamation.\footnote{Le $P$ de cet
énoncé est écrit trois fois par-dessus un $M$, et dans le noyau ce qui précède
$M \otimes P$ est biffé et surchargé ; la transcription retient la lecture
corrigée, et c'est elle qui est reprise ici.}

\emph{La réponse à 19.5.6 (iii), et quatre remarques.}

« L'implication b) $\Rightarrow$ c) de 19.5.5 est vraie sans qu'on doive faire
aucune hypothèse de séparation. Pour le voir, il suffit de lire le cas $n = 1$
de la démonstration “à la Schwartz”. » Il joint, vu l'urgence signalée, une
rédaction de la démonstration, « copiée sur la primitive » — cette rédaction
n'est pas dans le dossier.

Puis, en quatre tirets, ce qui n'appartient à aucun théorème :

\begin{itemize}
  \item « Je n'ai pas encore compris 19.7. »
  \item « J'ai beaucoup aimé le projet EGA I § 0 jusqu'à la pg 10, où la
  topologie du spectre tombe du ciel. Ne faudrait-il pas parler de topologie
  avant de parler d'espace topologique ? »
  \item « Je ne suis pas capable de rédiger une théorie naïve utilisable des
  intersections »,
  \item et rien d'autre à dire sur le « plan des EGA » pour le moment.
\end{itemize}

\pagerange{55}{56}
\subsection*{c) La lettre du 15 décembre 1966 (pages 55 et 56)}

Deux pages, et la dernière des trois. Elle répond à une lettre du 28 — qui n'est
pas dans le dossier — et elle revient sur l'objection d'octobre : « Après avoir
tourné 7 fois ma langue dans ma bouche, je répète mon objection à l'hypothèse
“tout quotient de tout sous-module de $M$ est séparé pour la topologie
$J$-adique”, qui \emph{n'est pas} en général vérifiée pour $M$ de type fini sur
la $A$-algèbre $B$, $A$ et $B$ noethériens, $JB \subset R(B)$. »

Le contre-exemple est en trois lignes. $A$ noethérien local, $J = \mathfrak{m}$,
$B = A[[T]]$, $M = B$. Comme $A$-module, $B \simeq A^{\mathbf{N}}$, donc tout
module quotient de $A^{(\mathbf{N})}$ est quotient d'un sous-$A$-module de $B$ ;
et « un module ayant un système dénombrable de générateurs n'a aucune raison
d'être séparé ! ».\footnote{L'indice $s$ de $A_{s}^{\mathbf{N}}$ et
$A_{s}^{(\mathbf{N})}$ est écrit ainsi sur le feuillet et n'est pas expliqué
dans la lettre ; il est omis ici, l'argument n'en ayant pas besoin. Deligne
concède immédiatement la limite de son objection : « Bien sûr, les
démonstrations n'utilisent que des quotients et sous-modules comme
$B$-modules. » L'objection porte donc sur le libellé de l'hypothèse et non sur
la validité des démonstrations qui l'invoquent.}

Deux post-scriptum. Le premier donne son adresse militaire en Allemagne. Le
second est le plus lourd :

\begin{quote}
Je commence à lire IV § 16. La démonstration de 16.1.5 est canulée, comme on le
voit en prenant $X = \mathrm{Spec}(k[T])$, $Y = \mathrm{Spec}(k[T]) \sqcup
\mathrm{Spec}(k)$, tous deux affins.
\end{quote}

Avec les notations de 16.1.15 (i), $J = 0$, mais les voisinages infinitésimaux
de $Y$ ne sont pas triviaux en $\mathrm{Spec}(k)$. Il ajoute : « J'ignore si la
proposition est vraie. »\footnote{La marge gauche porte un croquis de Deligne :
deux accolades « $Y$ » et « $X$ » devant deux traits horizontaux superposés,
celui de $Y$ portant un point isolé au-dessus, et une flèche verticale
descendante de $Y$ vers $X$ — la droite affine avec un point disjoint,
au-dessus de la droite affine. Le contre-exemple est donc un dessin d'une
ligne, et c'est un dessin qui suffit.}

Et enfin : « même en admettant 16.1.5, 16.1.6 n'est pas convaincant en général :
si $B$ est de présentation finie sur $A$, j'ignore si le noyau de $A \to B$ est
un idéal de type fini. C'est sans doute faux. »

C'est donc le paragraphe 16 — celui-là même que la fiche de la page 51 demande
de vérifier, celui dont le tapuscrit des pages 52 et 53 va retirer une
définition — qui reçoit ici, deux mois plus tard, son second coup.

\pagerange{66}{80}
\section*{XII. Divers (pages 66 à 80)}

Une feuille kraft par ailleurs vierge porte, en haut à droite et suivi de points
de suspension, le seul mot « Divers …. ». Ce qu'elle couvre est quinze feuillets
sans lien entre eux, qui sont des compléments au chapitre IV rangés dans aucun
ordre. On les prend comme ils viennent.

\subsection*{La structure d'une démonstration du Main Theorem (pages 66 et 67)}

Ce n'est pas une démonstration : c'est le \emph{plan} d'une démonstration, et
c'est à ce titre qu'il vaut d'être lu, car on y voit comment une preuve se
décompose en couches d'hypothèses qu'on enlève une à une.

\begin{enumerate}
  \item[a)] Démonstration par voie globale, dans le cas $Y$ localement
  noethérien et $f$ quasi-projectif.
  \item[b)] Extension au cas où $f$ est quasi-fini séparé de présentation finie
  et $Y$ quasi-compact quasi-séparé. On prouve que la question est alors locale
  sur $Y$ (IV 8.12.3), ce qui ramène au cas noethérien ; et que l'hypothèse
  faite implique $f$ quasi-affine (IV 8.11.2), ce qui ramène aux conditions de
  a).
  \item[c)] Élimination de l'hypothèse « $f$ de présentation finie », ne
  gardant que « $f$ quasi-fini séparé ».
\end{enumerate}

L'\emph{élimination de la méthode globale} est traitée à part : on se ramène au
cas où $Y$ est local essentiellement de type fini sur $\mathrm{Spec}\,\mathbf{Z}$,
on passe aux complétés et l'on fait une récurrence sur la dimension, en
utilisant que le passage au complété est compatible avec la fermeture intégrale
(IV 6.14.4).

Et pour c), deux voies : ou bien supposer $f$ quasi-projectif et passer à la
limite dans $X = \varprojlim X_{i}$ ; ou bien prouver que $f$ est nécessairement
quasi-affine, donc quasi-projectif (IV 18.12.12), par la technique de
localisation étale. Il note alors l'objection qui pèse sur la seconde voie, et
c'est l'observation la plus intéressante du feuillet : \emph{la technique de
localisation étale utilise elle-même le Main Theorem}, sous la forme locale en
haut qu'établit b). Le choix reste ouvert « entre la démonstration chronologique
globale, ou démonstration purement locale… », et le feuillet s'arrête là.

\subsection*{Morphismes essentiellement de présentation finie, et essentiellement affines (pages 69 et 70)}

Deux définitions et deux propositions, sous ses propres titres.

Soit $f : X \to Y$ un morphisme. On considère les conditions :

\begin{enumerate}
  \item[(i)] $f$ est la limite projective d'un système projectif filtrant
  essentiellement affine de $Y$-préschémas de présentation finie ;
  \item[(ii)] $f$ se factorise en $X \xrightarrow{f_{1}} X_{0}
  \xrightarrow{f_{0}} Y$ avec $f_{0}$ de présentation finie et $f_{1}$ affine ;
  \item[(iii)] comme (ii), mais avec $X = \mathrm{Spec}(\mathcal{A})$ où
  $\mathcal{A}$ est une Algèbre quasi-cohérente sur $X_{0}$ limite inductive
  filtrante d'Algèbres de présentation finie.
\end{enumerate}

Alors (i) $\Leftrightarrow$ (iii) $\Rightarrow$ (ii), et tout est équivalent si
$Y$ est localement noethérien.

\emph{Définition.} $f$ est \emph{essentiellement de présentation finie} s'il
satisfait (ii) ; $X$ l'est si $X \to \mathrm{Spec}\,\mathbf{Z}$ l'est. On note
que cela entraîne $f$ quasi-compact quasi-séparé.

La proposition qui suit donne la boîte à outils habituelle : la présentation
finie implique la notion ; un morphisme affine de but localement noethérien, ou
de but quasi-compact quasi-séparé avec faisceau inversible ample, l'est ; une
immersion de but quasi-compact quasi-séparé l'est ; la notion est stable par
changement de base et par composition ; et si $g$ est essentiellement de
présentation finie (avec $S$ quasi-compact quasi-séparé, ou $g$ de présentation
finie), alors $gf$ l'étant, $f$ l'est.

Une seconde proposition, plus fine : si $S$ est universellement japonais — par
exemple $S = \mathrm{Spec}\,\mathbf{Z}$ — et si $X$ est essentiellement de
présentation finie sur $S$, alors pour que $X$ soit géométriquement unibranche
il faut et il suffit qu'il soit de la forme $\varprojlim X_{i}$ avec les $X_{i}$
géométriquement unibranches (resp. normaux).

Le feuillet 70 reprend la seule notion « essentiellement affine » — même
factorisation, sans la condition de présentation finie sur $f_{0}$ — et
commence une liste de propriétés dont le (i) est écrit et le (ii) resté vide.
Le feuillet s'arrête là.

\subsection*{Les préschémas dont les anneaux locaux réduits sont des corps (page 72)}

Un feuillet de conditions équivalentes, et c'est une caractérisation des
préschémas de dimension zéro. La liste, dans sa numérotation finale :

\begin{enumerate}
  \item[(i)] Tout point est fermé.
  \item[(ii)] Tout point admet un voisinage ouvert localement compact ; (ii b)
  idem, compact ; (ii ter) tout ouvert affine est compact ; (ii quater) si $X$
  est quasi-séparé, $X$ est localement compact.
  \item[(iii)] Les anneaux locaux de $X_{\mathrm{réd}}$ sont des corps.
  \item[(iv)] $X^{\mathrm{cons}} \to X$ est un homéomorphisme (ou : un
  isomorphisme) ; (iv bis) toute partie localement constructible est ouverte ;
  (iv ter) si $X$ est quasi-séparé, toute partie constructible est ouverte ;
  (iv quater) tout ouvert affine est fermé.
\end{enumerate}

$X^{\mathrm{cons}}$ désigne l'espace $X$ muni de la topologie constructible.
L'équivalence de (i) et (iii) est triviale — tout point fermé signifie tout
idéal premier maximal, c'est-à-dire dimension zéro. Les autres disent que
l'espace est un espace localement compact totalement discontinu.

Le feuillet porte ensuite une chaîne d'implications, et il faut la lire avec un
avertissement. \emph{La liste a été renumérotée de $(v\ldots)$ en $(iv\ldots)$
et la chaîne ne l'a pas été.}\footnote{La transcription note que les numéros
(iv bis), (iv ter), (iv quater) sont récrits en surcharge sur « (v bis) »,
« (v ter) », « (v quater) », et que les renvois gardent la numérotation en (v).
Les lignes « (v) $\Leftrightarrow$ (v bis) », « (v bis) $\Rightarrow$ (v ter)
$\Leftarrow$ (v quater) », « (i) $\Rightarrow$ (v quater) » se lisent alors sans
peine en remplaçant chaque $v$ par le $iv$ correspondant. Une ligne reste sans
référent : « (iv) $\Leftrightarrow$ (v) », qui, après renumérotation, opposerait
(iv) à lui-même. Elle est reportée ici telle quelle, sans être interprétée : la
liste d'origine comportait vraisemblablement un membre de plus, supprimé au
moment de la renumérotation, et il n'est plus possible de dire lequel.} Les
justifications qu'il donne sont d'un mot — « if », pour « immédiat » ou « il
faut » — sauf une, (i) $\Rightarrow$ (iv quater) : « car l'adhérence est formée
des spécialisations ».

Le feuillet ajoute enfin deux choses. Un N.B. : tout espace localement compact
totalement discontinu muni d'un faisceau d'anneaux dont les fibres sont des
corps est un préschéma. Et un corollaire, qui est le véritable intérêt de la
page :

\emph{Corollaire.} Conditions équivalentes : (i) $X$ est « localement compact »
(pas nécessairement séparé) et réduit ; (ii) tout Module sur $X$ est plat ;
(ii bis) tout Module de présentation finie sur un ouvert affine de $X$ est
plat ; (ii ter) si $X$ est quasi-séparé, tout Module quasi-cohérent sur $X$ est
plat ; (iii) $X^{\mathrm{cons}} \to X$ est un isomorphisme d'espaces annelés.

C'est la caractérisation des anneaux \emph{absolument plats}, dits aujourd'hui
von Neumann réguliers : un anneau sur lequel tout module est plat. La condition
(iii) est l'énoncé géométrique correspondant — le spectre coïncide avec son
spectre constructible.

\subsection*{Deux feuillets de tapuscrit (pages 73 et 78)}

Le premier, paginé 41, énonce un théorème de descente effective pour les
faisceaux étales. Ses annotations manuscrites sont des consignes de mise en page
pour la dactylographe et ne concernent pas le contenu.

\emph{Théorème 9.4.} Soit $f : X' \to X$ un morphisme surjectif de préschémas.
On suppose vérifiée l'une des hypothèses : a) $f$ entier ; b) $f$ propre ; c)
$f$ plat et localement de présentation finie ; d) $X$ discret (par exemple le
spectre d'un corps). Alors $f$ est un \emph{morphisme de descente effective}
pour la catégorie fibrée des faisceaux étales : le foncteur $f^{*}$ induit une
équivalence entre la catégorie des faisceaux sur $X$ et celle des faisceaux sur
$X'$ munis d'une donnée de descente relativement à $f$.

La démonstration commencée (9.4.1) définit, à partir d'un faisceau $F'$ muni
d'une donnée de descente, le foncteur $G(Y) = \mathrm{Ker}\bigl(F'(Y')
\rightrightarrows F''(Y'')\bigr)$, constate que c'est un faisceau pour la
topologie étale, et construit un homomorphisme canonique injectif $G \to
f_{*}(F')$. Le feuillet s'interrompt là.

Le second tapuscrit, paginé 24 et corrigé de sa main, est la démonstration d'un
énoncé (i) par récurrence noethérienne. On réduit à montrer que $f_{*}F'$ est
localement constant et constructible sur un ouvert non vide convenable ; on
prend $x$ un point maximal de $X$ ; en remplaçant $X$ par un revêtement étale
d'un voisinage de $x$, on rend radiciel chaque point de $X'$ au-dessus de $x$,
de sorte que $X' \to X$ est somme de morphismes radiciels au voisinage de $x$
(EGA IV.9.6.1) ; en rétrécissant encore on rend $F'$ constant, et alors
$f_{*}F'$ est produit d'un nombre fini de faisceaux constructibles constants,
donc constructible et constant.\footnote{Les corrections manuscrites sont
portées dans le texte et signalées par la transcription : « ne » ajouté à
« récurrence noethérien », « convenable » substitué à « quelconque » biffé,
« (2.4 (i)) » écrit à la main dans un blanc du tapuscrit, « EGA » ajouté devant
la référence 9.6.1, et « Rétrécissant » substitué à « Remplaçant ». Toutes sont
reprises ici.}

\subsection*{La condition $(h_{0})$ (pages 74 à 76)}

\emph{Définition.} Un morphisme $f : T \to S$ vérifie $(h_{0})$ lorsque
l'application $\mathrm{OF}(S) \to \mathrm{OF}(T)$ est injective, ce qui équivaut
à dire que pour toute partie ouverte et fermée non vide $\mathcal{U}$ de $S$,
$f^{-1}(\mathcal{U})$ est non vide ; et encore à dire que le foncteur $S'
\mapsto T' = T \times_{S} S'$ des revêtements de $S$ aux revêtements étales de
$T$ est \emph{conservatif}. Cela entraîne qu'il est pleinement fidèle.

Deux cas immédiats : si $S$ est connexe non vide, la condition équivaut à $T
\neq \emptyset$ ; et si $S = \coprod_{i} S_{i}$ avec les $S_{i}$ connexes non
vides, elle équivaut à $f^{-1}(S_{i}) \neq \emptyset$ pour tout $i$. C'est donc
exactement la condition de surjectivité sur les composantes connexes.

\emph{Proposition.} Soient $f : T \to S$ un morphisme de $k$-préschémas ($k$ un
corps), $R'$ un $k$-préschéma, et $f' : T' \to S'$ le morphisme déduit par
changement de base. Si $f$ satisfait $(h_{0})$, alors $f'$ aussi ; et
réciproquement si $R' \neq \emptyset$.

Le sens facile est celui de la réciproque : $S' \to S$ étant surjectif, une
partie ouverte et fermée de $S$ dont l'image inverse dans $S'$ est vide est
vide. Le sens direct se traite fibre par fibre, ce qui ramène à $R' =
\mathrm{Spec}\,k'$ avec $k'$ un corps, que l'on peut supposer algébriquement
clos, puis clôture algébrique de $k$ ; on conclut alors par un argument galoisien
— l'image directe d'une partie ouverte et fermée de $S'$ est stable par
$\mathrm{Gal}(k'/k)$, donc ouverte et fermée dans $S$, et son image inverse dans
$T$ étant vide, elle est vide.\footnote{Les quatre premières lignes de la page
75, qui portent cet argument, sont un palimpseste : deux lignes biffées, deux
ajouts interlinéaires, et le $\mathcal{U}$ qui nomme l'image directe écrit en
exposant au-dessus du mot « (fermée) ». L'ordre des membres donné par la
transcription est celui qu'elle a cru lire, et la reconstitution ci-dessus suit
le sien. Un lecteur ayant le feuillet devrait y revenir : c'est, avec la page
77, l'endroit le moins sûr du quatrième lot.} Le feuillet note en marge :
« Cf simplification éventuelles de EGA IV 4 ».

\emph{Corollaire 1.} Si $f : T \to S$ satisfait $(h_{0})$ et si $R' \neq
\emptyset$, le carré

\begin{tikzcd}
  \mathrm{OF}(S) \arrow[r, hook] \arrow[d] & \mathrm{OF}(T) \arrow[d] \\
  \mathrm{OF}(S') \arrow[r, hook] & \mathrm{OF}(T')
\end{tikzcd}

est cartésien. L'injectivité de $\mathrm{OF}(S) \to \mathrm{OF}(T)$ donne
l'injectivité de $\mathrm{OF}(S) \to \mathrm{OF}(T) \times_{\mathrm{OF}(T')}
\mathrm{OF}(S')$, et la surjectivité est immédiate par descente, en utilisant
que $\mathrm{OF}(S'') \to \mathrm{OF}(T'')$ est injective.

Le feuillet 76 est une page de croquis, sans titre : trois diagrammes et
quelques lignes, à propos du prolongement d'un morphisme $A_{T} \to B_{T}$ en
un morphisme $A \to B$.\footnote{Plusieurs traits de ces diagrammes sont sans
pointe sur le feuillet et un croquis est redessiné à partir d'un dessin où deux
couples $(A_{0}, B_{0})$ et $(A, B)$ sont chacun reliés à deux sommets par des
traits obliques. Ce qui est reproduit dans la transcription est une
reconstruction, et n'est pas repris ici comme diagramme.} Ce qui s'en dégage est
un énoncé d'unicité et un critère d'existence : sous les deux hypothèses que
toute partie ouverte et fermée de $S''$ d'image inverse vide dans $T''$ est
vide — c'est-à-dire $(h_{0})$ pour $T'' \to S''$ — et que $S' \to S$ est un
morphisme de descente pour la catégorie des schémas relatifs, il existe
\emph{au plus} un $u : A \to B$ induisant un $v : A_{T} \to B_{T}$ donné, et il
en existe un si et seulement s'il en existe un après le changement de base $S'
\to S$.

\subsection*{Les conditions sur $Y \to X$ : composantes connexes et revêtements étales (pages 77 et 79)}

Deux feuillets au crayon, sur le dos de deux tapuscrits qui transparaissent, et
qui portent \emph{deux fois la même liste} sous deux formes. Ce sont les plus
difficiles du dossier, et la comparaison des deux rédactions est le seul moyen
de les lire.

Rappelons la clef : chaque condition porte sa forme forte entre parenthèses, et
la liste vaut donc deux fois.

\emph{Conditions équivalentes pour $f : Y \to X$.}

\begin{enumerate}
  \item[(i)] Pour tout $X'$ étale sur $X$, posant $Y' = Y \times_{X} X'$,
  l'application $\mathrm{OF}(X') \to \mathrm{OF}(Y')$ est injective
  (bijective).
  \item[(i bis)] ($X$, $Y$ localement noethériens) Pour tout tel $X'$,
  $\pi_{0}(Y') \to \pi_{0}(X')$ est surjective (bijective).
  \item[(ii)] $\Gamma(X'/X) \to \Gamma(Y'/Y)$ est injective (bijective).
  \item[(iii)] Le foncteur $X' \mapsto Y' = X' \times_{X} Y$, de la catégorie
  des revêtements étales finis de $X$ dans celle des revêtements étales finis
  de $Y$, est fidèle (pleinement fidèle).
\end{enumerate}

C'est la condition de type Lefschetz sur le groupe fondamental : dans la lecture
forte, la pleine fidélité du foncteur des revêtements étales finis est
exactement la surjectivité de $\pi_{1}(Y) \to
\pi_{1}(X)$.\footnote{L'identification est celle de la théorie de Galois
grothendieckienne, et n'est pas faite sur les feuillets : ils énoncent les
quatre conditions et les relient, sans nommer le groupe fondamental. Elle est
donnée ici comme le nom moderne de ce que les pages construisent.}

Le feuillet 77 porte, entre les conditions, les raisons des équivalences,
écrites en deux flèches courbes dans la marge : de (iii) vers (i), « considération
des graphes » ; de (i) vers (iii), « par considération des Hom » ; et sous elles
la ligne qui fait le pont, « Car $\mathrm{OF}(X') \simeq \Gamma(I_{X'}/X')$ ».

Le feuillet 79 ajoute une quatrième condition — $f(Y)$ \emph{dense} dans $X$ —
et, surtout, la \emph{réserve} que la page 77 ne portait pas. L'implication (i)
$\Rightarrow$ (ii) ne vaut « comme avant » que tant qu'on se borne aux schémas
\emph{séparés} sur $X$, ou qu'il ne s'agit que de l'énoncé faible ; sinon « il
semble qu'il faille supposer $f$ universellement [quelque chose] et [qu'il]
reste tel après changement de base ». Il note enfin que si $X$ et $Y$ sont
essentiellement de présentation finie et qu'on impose des conditions stables par
tout changement de base fini, « alors c'est de nouveau OK. (cf
SGAA) ».\footnote{Deux mots de cette réserve n'ont pas pu être lus, dont celui
qui suit « universellement » — c'est-à-dire précisément la condition à imposer.
Le crochet à droite de (ii), sur la page 79, porte la mention « avec réserve »,
qui est ce qui a conduit à la chercher. Rien n'est reconstitué ici : la réserve
est signalée, son contenu manque.} Le haut du feuillet 77, au-dessus d'un trait
horizontal, porte des croquis étrangers à la liste — un $X'$ biffé d'où partent
deux flèches vers $Y$, et trois modules de différentielles dont un biffé.

\subsection*{Croquis : images directes et changement de base (page 80)}

Un feuillet de croquis à l'encre, tracés en travers d'un tapuscrit filmé
tête-bêche qui n'est pas de sa main. Seuls des fragments portent :
\[
  i_{*}(\Lambda_{V}) = \Lambda_{C} , \qquad
  R^{1} i_{*}(\Lambda_{V}) = \mathcal{H}^{-1}_{Z,C} ,
\]
ce dernier souligné deux fois et suivi d'un renvoi « 2.14. » ; puis deux carrés
cartésiens et la flèche de changement de base
\[
  u^{*}\bigl(R^{i} p_{*}(\overline{F})\bigr) \longrightarrow
  R^{i} p_{Y*}\bigl(u'^{*}(\overline{F})\bigr) .
\]
C'est le morphisme de changement de base pour les images directes supérieures,
noté pour lui-même. Le feuillet ne dit pas sous quelles hypothèses il serait un
isomorphisme, et rien n'est supposé ici à sa place.

\pagerange{82}{89}
\section*{XIII. IV 5 : dimension, une $\partial$-catégorie, la clôture intégrale (pages 82 à 89)}

Une chemise orange porte, à l'encre puis au crayon : « IV 5 », un mot biffé qui
ne se lit plus, « Compléments d'ordre divers ».

\subsection*{Un lemme de codimension, et deux corollaires (pages 82 et 83)}

\emph{Lemme 1.} Soient $X$ un préschéma noethérien, $Y$ une partie fermée non
vide, $(Y_{i})_{i \in I}$ l'ensemble de ses composantes irréductibles,
$(X_{k})_{k \in K}$ celles de $X$. Posons $r = \mathrm{Codim}(Y, X)$, $I'$
l'ensemble des $i$ tels que $\mathrm{Codim}(Y_{i}, X) = r$, et pour $i \in I'$,
$K_{i}$ l'ensemble des $k$ tels que $X_{k} \supset Y_{i}$ et
$\mathrm{Codim}(Y_{i}, X_{k}) = r$. Soit enfin $x$ une section d'un Module
inversible $L$ sur $X$ telle que $V(x) \supset Y$. Alors :

\begin{enumerate}
  \item[1\textsuperscript{o})] $\mathrm{Codim}(Y, V(x)) \geqslant r - 1$, avec
  égalité si et seulement s'il existe $i \in I'$ et $k \in K_{i}$ tels que
  $V(x) \not\supset X_{k}$ ;
  \item[2\textsuperscript{o})] si $X$ est local et $r > 0$, il existe $x \in
  \mathfrak{m}$ tel que $\mathrm{Codim}(Y, V(x)) = r - 1$ et $\dim V(x) = \dim X
  - 1$.
\end{enumerate}

La démonstration de 1\textsuperscript{o} est l'enchaînement du feuillet, et
c'est le théorème de l'idéal principal de Krull qui la porte :
$\mathrm{Codim}(Y, V(x)) = \mathrm{Inf}_{i}\, \mathrm{Codim}(Y_{i}, V(x))$, et
\[
  \mathrm{Codim}(Y_{i}, V(x)) = \dim \mathcal{O}_{V(x), y_{i}}
  \geqslant \dim \mathcal{O}_{X, y_{i}} - 1 = \mathrm{Codim}(Y_{i}, X) - 1
  \geqslant \mathrm{Codim}(Y, X) - 1 .
\]
L'égalité impose donc les deux égalités simultanées
$\mathrm{Codim}(Y_{i}, V(x)) = \mathrm{Codim}(Y_{i}, X) - 1$ et
$\mathrm{Codim}(Y_{i}, X) = \mathrm{Codim}(Y, X)$, c'est-à-dire $i \in I'$ et
l'existence d'un $k \in K_{i}$ avec $V(x) \not\supset X_{k}$ — autrement dit
\emph{la section doit couper effectivement au moins une composante de la bonne
codimension}.

Pour 2\textsuperscript{o} : comme $r > 0$, aucune composante irréductible
$X_{\alpha}$ de $X$ n'est contenue dans $Y$, donc $\mathfrak{m} \not\subset
\mathfrak{p}_{\alpha}$ pour tout $\alpha$, et le lemme d'évitement fournit $x
\in \mathfrak{m}$ n'appartenant à aucun des $\mathfrak{p}_{\alpha}$.\footnote{Le
feuillet 83 porte, entièrement barré de traits diagonaux, une première rédaction
de 2\textsuperscript{o} beaucoup plus longue, qui essayait de choisir $x$
composante par composante ; elle est remplacée par les trois lignes données ici.
En marge, à sa hauteur, une phrase dont ne se lisent que les mots « Je
modifie », « conditions » et « indiquées ».}

\emph{Corollaire 1.} Soient $A$ un anneau noethérien et $\mathfrak{m}$ un idéal
de hauteur $r$. Alors il existe $x_{1},\ldots,x_{r} \in \mathfrak{m}$ avec
$\mathrm{ht}\bigl(\sum x_{i}A\bigr) = r$ ; si $\mathfrak{m}$ est premier, il est
minimal au-dessus de $(x_{1},\ldots,x_{r})$. Si $A$ est local et $\mathfrak{m}$
son idéal maximal, les $x_{i}$ forment un système de
paramètres.\footnote{L'énoncé du feuillet hésite entre un anneau local et un
anneau noethérien muni d'un idéal $\mathfrak{m}$ : « local » n'est pas biffé
dans la première ligne, mais la clause finale « si $A$ local » n'aurait aucun
sens si $A$ l'était déjà. C'est donc la seconde lecture qui est retenue, et
c'est elle qui rend la dernière clause informative. Le feuillet écrit la
conclusion « $\mathrm{ht}(\mathfrak{m}/(x_{1},\ldots,x_{r})) = 0$ », ce qui est
la minimalité de $\mathfrak{m}$ au-dessus de l'idéal engendré ; elle ne vaut
telle quelle que pour $\mathfrak{m}$ premier, et c'est sous cette forme qu'elle
est donnée ici. La ligne suivante du feuillet, « Prenons $\mathfrak{m}$
premier », confirme que c'est vers ce cas qu'il va.}

\emph{Corollaire 2.} Soient $A$ local noethérien, $\mathfrak{p}$ un idéal
premier et $r = \dim A_{\mathfrak{p}}$. Alors il existe $x_{1},\ldots,x_{r} \in
\mathfrak{p}$ définissant un système de paramètres de $A_{\mathfrak{p}}$ ; et si
$\dim A/\mathfrak{p} + \dim A_{\mathfrak{p}} = \dim A$ — par exemple si $A$ est
caténaire et équidimensionnel — ils font partie d'un système de paramètres de
$A$.\footnote{Deux départs ici. Le feuillet écrit « font partie d'un système de
paramètres de $A/\mathfrak{p}$ », et la transcription note qu'on attendrait
$A$ : c'est bien $A$, car les $x_{i}$ appartiennent à $\mathfrak{p}$ et sont
donc \emph{nuls} dans $A/\mathfrak{p}$, où ils ne peuvent faire partie de rien.
Et la complétion en un système de paramètres de $A$ n'est pas automatique : elle
demande $\dim A/(x_{1},\ldots,x_{r}) = \dim A - r$, c'est-à-dire la formule de
dimension rappelée ici. Le feuillet ne la mentionne pas ; elle est suppléée.}

La marge gauche du feuillet 83, écrite de bas en haut sur toute la hauteur,
ajoute un corollaire de plus : pour $x_{1},\ldots,x_{p} \in \mathfrak{m}_{A}$,
$A$ local noethérien, les conditions $\mathrm{ht}(\sum A x_{i}) = p$ — l'autre
inégalité étant toujours vraie — et « $x_{1}$ ne s'annule sur aucune composante
irréductible de $\mathrm{Spec}(A)$, …, $x_{n}$ ne s'annule sur aucune composante
irréductible de $\mathrm{Spec}\,A/\sum_{i<n} x_{i}A$ » sont
équivalentes.\footnote{La fin de la condition (ii), en très petits caractères,
est lue d'après le sens par la transcription : le quotient et son indice de
sommation y sont donnés comme incertains. La forme donnée ici est celle qui rend
l'équivalence vraie, et c'est la définition usuelle d'un système de paramètres
partiel.}

\subsection*{Une $\partial$-catégorie (page 84)}

Un feuillet entier, barré de traits diagonaux mais parfaitement lisible sous la
rature, qui pose les premiers axiomes d'une notion de catégorie graduée à
translation.

\emph{Définition.} Une $\partial$-catégorie $\mathcal{H}$ est une catégorie
additive \emph{graduée} — les $\mathrm{Hom}(X,Y)$ sont des groupes gradués et la
composition est additive sur les degrés — satisfaisant les axiomes qui suivent.
On note $\mathcal{H}_{(0)}$ la catégorie où l'on ne garde que les morphismes de
degré $0$.

\emph{Axiome $\partial 1$.} Pour tout objet $X$, il existe un objet $X(1)$ et
un isomorphisme fonctoriel en $X$ sur $\mathcal{H}_{(0)}$
\[
  \mathrm{Hom}_{1}(X, Y) \xrightarrow{\;\sim\;} \mathrm{Hom}_{0}(X, Y(1)) .
\]

Le feuillet donne aussitôt la forme concrète de cet axiome, qui lui est
équivalente : il existe un objet $X(1)$ et un isomorphisme $\alpha_{X} : X(1)
\to X$ \emph{de degré $1$}, et un objet $X(-1)$ et un isomorphisme $\beta_{X} :
X(-1) \to X$ de degré $-1$.\footnote{Les deux formulations sont écrites l'une
après l'autre sur le feuillet, la seconde introduite par un « il existe »
interlinéaire : c'est une reprise, non un second axiome. Le passage de la
première à la seconde est la représentabilité rendue explicite — $Y(1)$
représente $\mathrm{Hom}_{1}(-, Y)$, et le morphisme universel est l'élément
$i_{Y} \in \mathrm{Hom}_{1}(Y(1), Y)$ que le feuillet nomme, en notant que
$(Y(1), i_{Y})$ est déterminé à isomorphisme près.}

On en déduit, par composition avec $\alpha_{X}$ et $\beta_{X}$, quatre
isomorphismes pour tout $Y$ :
\[
  \text{I}
  \begin{cases}
    \mathrm{Hom}_{n}(X, Y) \xrightarrow{\sim} \mathrm{Hom}_{n+1}(X(1), Y) \\
    \mathrm{Hom}_{n}(Y, X(1)) \xrightarrow{\sim} \mathrm{Hom}_{n+1}(Y, X)
  \end{cases}
  \qquad
  \text{II}
  \begin{cases}
    \mathrm{Hom}_{n}(X, Y) \xrightarrow{\sim} \mathrm{Hom}_{n-1}(X(-1), Y) \\
    \mathrm{Hom}_{n}(Y, X(-1)) \xrightarrow{\sim} \mathrm{Hom}_{n-1}(Y, X)
  \end{cases}
\]
et la marge ajoute, contre la première ligne de I, l'identification
\[
  \mathrm{Hom}_{n}(X,Y) = \mathrm{Hom}_{n-1}(X, Y(1)) .
\]

Le feuillet s'arrête après ce premier axiome : le second n'est pas dans le
dossier, et la notion n'est donc pas définie. Ce qu'elle serait devenue se
reconnaît toutefois sans peine — $X \mapsto X(1)$ est un foncteur de
translation, et une catégorie additive graduée munie d'une translation est le
cadre dans lequel se formulent les catégories triangulées et, plus tard, les
catégories différentielles graduées.\footnote{La rature en diagonale qui barre
le feuillet est de lui. Ce rapprochement est le nôtre et ne prétend qu'à
situer : rien sur la page ne parle de triangles, de cônes, ni d'axiome
octaédral, et le seul axiome présent est celui de la translation.}

\subsection*{La finitude de la clôture intégrale (pages 86 à 88)}

Un argument écrit hors de l'ordre des pages. La \emph{Remarque} commence au
feuillet 88 et se poursuit au feuillet 87 ; la marge de 87 porte « cf verso », et
le verso est le feuillet 86, qui contient le lemme dont elle se sert. L'ordre de
lecture est donc 88, puis 87, puis 86.

\emph{Remarque (page 88).} Soient $A$ un anneau intègre \emph{quotient d'un
anneau régulier} et $K$ son corps des fractions. Pour que la clôture intégrale
de $A$ dans $K$ soit de type fini, il faut et il suffit que :

\begin{enumerate}
  \item[a)] l'ensemble des $\mathfrak{p}_{i} \in \mathrm{Spec}\,A$ tels que
  $A_{\mathfrak{p}_{i}}$ soit de dimension $1$ et non régulier soit \emph{fini} ;
  \item[b)] pour chacun d'eux, la normalisée de $A_{\mathfrak{p}_{i}}$ soit
  finie sur $A_{\mathfrak{p}_{i}}$.
\end{enumerate}

La nécessité est évidente : le lieu non normal de $\mathrm{Spec}\,A$ est fermé,
et les $\mathfrak{p}_{i}$ en sont les points génériques de codimension $1$.

La suffisance est la construction que le reste des trois feuillets exécute.
Comme $A$ est quotient d'un régulier, le lieu où $A_{x}$ satisfait $(S_{2})$ est
\emph{ouvert}, et contient un ouvert $X_{f}$ ; grâce à a) et b), il existe un
sous-anneau $B$ de $K$ contenant $A$, fini sur $A$, tel que les
$B_{\mathfrak{p}_{i}}$ soient normaux. Posons $Z = \mathrm{Supp}(B/A)$ : il est
contenu dans $X - X_{f}$ et de codimension $\geqslant 2$. Soient $U = X - Z$,
$i : U \to X$ l'immersion, $\mathcal{B} = \widetilde{B}$, et
\[
  \mathcal{C} = i_{*} i^{*}(\mathcal{B}) .
\]
C'est le prolongement de $\mathcal{B}$ à travers $Z$, et comme $Z$ est de
codimension $\geqslant 2$ et $\mathcal{B}$ satisfait $(S_{2})$, le faisceau
$\mathcal{C}$ est \emph{cohérent} — c'est le théorème de prolongement à travers
une codimension $2$.

On montre ensuite que $\mathcal{C}$ est normal par le critère de Serre.

\emph{$(S_{2})$.} Par construction $\mathcal{C} \xrightarrow{\sim} i_{*}i^{*}
\mathcal{C}$, et sur $U$ on a $\mathcal{C} \simeq \mathcal{B}$, où
$\mathcal{B}_{x}$ est une algèbre finie et normale sur $\mathcal{O}_{X,x}$, donc
$(S_{2})$ comme anneau et par suite comme module ; $\mathcal{C}$ est donc
$(S_{2})$ partout, et $A \subset \mathcal{C} \subset K$.

\emph{$(R_{1})$.} Soit $\mathfrak{p}$ un idéal premier de $\mathcal{C}$ avec
$\dim \mathcal{C}_{\mathfrak{p}} = 1$, et $\mathfrak{q}$ l'idéal qu'il induit sur
$A$. Comme $A$ est quotient d'un régulier, $\dim A_{\mathfrak{q}} = 1$.\footnote{Le
feuillet porte cette raison entre crochets — « [$A$ est quotient d'un
régulier] » — et ne l'argumente pas. C'est là que l'hypothèse sert une seconde
fois : elle rend $A$ universellement caténaire, de sorte que les hauteurs se
conservent le long de l'extension finie birationnelle $A \subset \mathcal{C}$.
Le pas est suppléé ; la page ne le fait pas.} Alors ou bien
$A_{\mathfrak{q}}$ est normal, et $A_{\mathfrak{q}} = \mathcal{B}_{\mathfrak{p}}$
est normal ; ou bien il ne l'est pas, c'est-à-dire $\mathfrak{q}$ est l'un des
$\mathfrak{p}_{i}$, et alors $\mathcal{B}_{\mathfrak{p}}$ est normal par
construction, comme localisé de $B_{\mathfrak{q}}$. Dans les deux cas
$\mathcal{C}_{\mathfrak{p}}$ est régulier.

$\mathcal{C}$ est donc normal, fini sur $A$, et compris entre $A$ et $K$ : c'est
la clôture intégrale, et elle est finie.

\emph{Le lemme (page 86).} Soient $A$ un anneau, $B$ une algèbre finie sur $A$,
$M$ un $B$-module de type fini. On peut supposer $A \subset B$ et
$\mathrm{Supp}\,M = \mathrm{Spec}\,B$.

\begin{enumerate}
  \item[(i)] Si $M$, vu comme $A$-module, est $(S_{k})$, alors $M$ est
  $(S_{k})$ sur $B$.
  \item[(ii)] Si, pour deux idéaux premiers $\mathfrak{q}$, $\mathfrak{q}'$ de
  $B$ au-dessus d'un même $\mathfrak{p}$ de $A$, on a $\dim M_{\mathfrak{q}} =
  \dim M_{\mathfrak{q}'}$, alors $M$ étant $(S_{k})$ sur $B$, il l'est sur $A$.
\end{enumerate}

La démonstration repose sur les deux formules qui comparent les invariants de
part et d'autre d'un morphisme fini : pour $\mathfrak{p}$ dans $A$,
\[
  \mathrm{prof}_{A_{\mathfrak{p}}} M_{\mathfrak{p}}
  = \mathrm{Inf}_{\mathfrak{q}\,|\,\mathfrak{p}}\;
    \mathrm{prof}_{B_{\mathfrak{q}}} M_{\mathfrak{q}} ,
  \qquad
  \dim_{A_{\mathfrak{p}}} M_{\mathfrak{p}}
  = \mathrm{Sup}_{\mathfrak{q}\,|\,\mathfrak{p}}\;
    \dim_{B_{\mathfrak{q}}} M_{\mathfrak{q}} ,
\]
la profondeur prenant l'\emph{inf} et la dimension le \emph{sup}. Pour (i), il
vient alors, pour $\mathfrak{q}$ au-dessus de $\mathfrak{p}$ :
\[
  \mathrm{prof}_{B_{\mathfrak{q}}} M_{\mathfrak{q}}
  \geqslant \mathrm{prof}_{A_{\mathfrak{p}}} M_{\mathfrak{p}}
  \geqslant \mathrm{Inf}\bigl(k,\, \dim_{A_{\mathfrak{p}}} M_{\mathfrak{p}}\bigr)
  \geqslant \mathrm{Inf}\bigl(k,\, \dim_{B_{\mathfrak{q}}} M_{\mathfrak{q}}\bigr) ,
\]
ce qui est $(S_{k})$ sur $B$. Pour (ii), l'hypothèse d'équidimensionnalité des
fibres donne $\dim_{A_{\mathfrak{p}}} M_{\mathfrak{p}} = \dim_{B_{\mathfrak{q}}}
M_{\mathfrak{q}}$ pour \emph{tout} $\mathfrak{q}$ au-dessus de $\mathfrak{p}$ —
le sup d'éléments égaux — et la profondeur est atteinte en un $\mathfrak{q}$
convenable, d'où la conclusion.\footnote{Sur le feuillet, les lettres
$\mathfrak{p}$ et $\mathfrak{q}$ \emph{changent de côté} entre l'énoncé et la
démonstration : dans l'énoncé $\mathfrak{q}$ est un premier de $B$ au-dessus de
$\mathfrak{p}$ dans $A$ ; dans la démonstration, à partir de la ligne biffée
« l'hypothèse implique… », c'est $\mathfrak{p}$ qui est dans $B$ et
$\mathfrak{q}$ dans $A$, la flèche $\mathfrak{p} \to \mathfrak{q}$ notant
« au-dessus de ». L'échange est sur le feuillet et la transcription ne l'a pas
régularisé, ce qui est sa règle ; il l'est ici, une fois pour toutes, dans le
sens de l'énoncé. Le bas du feuillet, une dizaine de lignes, est hachuré au
point de ne plus se lire ; seule surnage, en bas à gauche, la formule
« $\mathrm{prof} \geqslant \mathrm{Inf}(k, \dim)$ », qui est le ressort de tout
le lemme. Le bas du feuillet 87 porte une première rédaction, hachurée, du même
lemme, qui énonçait une \emph{équivalence} là où la version retenue ne donne que
deux implications sous hypothèses : c'est un affaiblissement délibéré, et il est
juste.}

\subsection*{Un brouillon, et une couche au crayon (page 89)}

Le feuillet 89 est un brouillon à l'encre pour le lemme précédent, barré de
traits diagonaux, où l'on retrouve les inégalités de profondeur et une remarque
qu'il souligne d'un trait ondulé : $B$ normal est $(S_{2})$ « comme anneau, mais
pas comme module ». C'est exactement la distinction dont la démonstration de la
page 87 avait besoin.

Sous l'encre, une couche antérieure au crayon, en anglais, écrite dans l'autre
sens du feuillet, et sans rapport avec le reste du dossier. Elle porte sur
$\varphi : \mathrm{Pic}(X') \to \mathrm{Pic}(Y')$ :

\begin{enumerate}
  \item[1)] si pour tout $x \in X' - Y'$ avec $\dim \overline{x} = 1$ on a
  $\mathrm{prof}\,A_{x} \geqslant 3$, alors $\varphi$ est injectif ;
  \item[2)] si de plus $\mathrm{prof}\,A_{x} \geqslant 4$ et
  $\mathrm{Pic}(\mathrm{Spec}\,A_{x} - \{x\}) = 0$, alors $\varphi$ est bijectif.
\end{enumerate}

Ce sont les conditions de type Lefschetz pour le groupe de Picard : la
profondeur en tout point du complémentaire mesure jusqu'à quel degré la
restriction à $Y'$ perd de l'information, $3$ suffisant pour l'injectivité et
$4$ pour la surjectivité.\footnote{Elles se lisent aujourd'hui comme les
conditions $\mathrm{Lef}$ et $\mathrm{Leff}$ du séminaire SGA 2 sur les
théorèmes de type Lefschetz. Le rapprochement est nôtre : la page ne nomme rien
et ne donne aucune démonstration. Les derniers fragments de la marge — une
intersection complète de dimension $\geqslant 4$, $\mathrm{Pic}(X - a) = 0$, un
idéal $\mathfrak{a} = \mathfrak{f}A + \mathfrak{g}A$ — sont au crayon, en partie
sous l'encre, et leur liaison avec 1) et 2) ne se lit pas.} C'est la seule page
du dossier qui soit écrite en anglais, et la seule qui regarde vers la
cohomologie.

\pagerange{91}{98}
\section*{XIV. IV 12 : un lieu non Cohen-Macaulay qui saute dans la mauvaise direction, et une proposition sur les algèbres graduées (pages 91 à 98)}

La dernière chemise orange porte « IV 12 » à l'encre et, au crayon, « Fibres des
morphismes plats de présentation finie ».\footnote{Le second chiffre est repassé
et pourrait aussi se lire $82$ ; le titre au crayon, qui est celui du
paragraphe 12 du chapitre IV, tranche en faveur de $12$.} Elle couvre huit
feuillets, en deux blocs que rien ne relie : un exemple au brouillon, et une
proposition mise au net.

\subsection*{L'exemple (pages 91 à 93)}

Un brouillon rapide, encre par-dessus crayon, avec un croquis à gauche. Ce
qu'il cherche est dit par la note qu'il écrit en accolade en haut à droite,
d'une encre plus noire que le titre :

\begin{quote}
Où le [lieu] des points non Coh. Mac. fait un saut dans la mauvaise
direction…
\end{quote}

L'objet est donc un \emph{contre-exemple}, et à ceci : on s'attendrait à ce que
la codimension du lieu non Cohen-Macaulay ne puisse pas augmenter quand on coupe
par une hypersurface. Elle le peut.

La construction est un recollement, et elle est décrite deux fois — abstraitement
à la page 91, puis explicitement à la page 93.

\emph{La forme générale (page 91).} On prend $X'$ un schéma local intègre de
dimension $n'$, de centre $x$, avec $Z'_{0}$ un point où il n'est pas
Cohen-Macaulay ; $X''$ un schéma local intègre de dimension $n'' > n'$ ; et $D$
un schéma local intègre de dimension $1$, muni de deux immersions fermées
\[
  u' : D \to X' , \qquad u'' : D \to X'' ,
  \qquad\text{avec}\qquad u'(D) \not\supset Z'_{0} ,
\]
cette dernière condition entourée et soulignée deux fois sur le feuillet. On
pose
\[
  X = X' \coprod_{D} X'' .
\]
On choisit alors $f_{0} \in \mathcal{O}_{D,x}$ non nul, on le relève en $f'$ et
$f''$, d'où $f \in \mathcal{O}_{X,x}$, et l'on a
\[
  V(f) \simeq V(f') \coprod_{V(f_{0})} V(f'') .
\]

\emph{Ce que l'exemple montre (page 92).} Si $V(f')$ et $V(f'')$ sont sans
idéaux immergés et si
\[
  V(f') \cap Z'_{0} = \{x\} , \qquad V(f'') \cap Z''_{0} \subset \{x\} ,
\]
alors $V(f)$ est sans idéaux immergés et est Cohen-Macaulay en dehors de
$\{x\}$. Le lieu où $X$ n'est pas Cohen-Macaulay contient $Z = Z'_{0} \cup
Z''_{0} \cup D$ — le recollement le long d'une courbe crée, tout le long de
$D$, des points de profondeur trop petite — et sa codimension dans $X$ est
majorée par $\mathrm{codim}_{X'} Z'_{0}$. Tandis que le lieu où $V(f)$ n'est pas
Cohen-Macaulay est réduit au point $x$, de codimension $\dim V(f) = n'' - 1
\geqslant n'$ dans $V(f)$. La codimension a donc \emph{augmenté} en passant de
$X$ à $V(f)$, et c'est le saut annoncé.\footnote{Le feuillet 92 est un brouillon
et sa chaîne d'inégalités ne se lit pas continûment : le membre « $X - V(f)$ est
de codim $\leqslant \mathrm{codim}_{X'} Z'_{0} \leqslant n'$ » porte un signe
surchargé qui pourrait être « $-1$ », plusieurs mots de liaison manquent, et
c'est la dernière inégalité, $n'' - 1 \geqslant n'$, qui est soulignée par le
sens. La formulation donnée ci-dessus est celle que la conjonction de la note
marginale et de ces fragments compose ; elle n'est pas une transcription de la
chaîne, et un lecteur qui voudrait s'en servir devrait refaire le calcul.}

\emph{La construction effective (page 93).} Pour $X'$, on prend le localisé au
sommet du cône projetant d'un schéma projectif normal $S'$ de dimension $d
\geqslant 3$ qui n'est \emph{pas} Cohen-Macaulay, plongé projectivement
normalement : alors $n' = d + 1 \geqslant 4$, et la condition $Z'_{0} \neq
\emptyset$, $Z'_{0} \neq \{x\}$ est vérifiée. Pour $X''$, le localisé à
l'origine d'un espace affine de dimension $n'' > n' = d+1$ — il est régulier,
donc $Z''_{0}$ est vide et la seconde condition d'intersection est
triviale. Pour $D$, $\mathrm{Spec}\,k[t]_{(t)}$, avec $u'$ obtenue en prenant
une générisation non générique de façon que $u'(D) \not\supset Z'_{0}$, et $u''$
une droite à l'origine. On choisit $f_{0}$ quelconque, puis $f'$ et $f''$ de
façon à vérifier $V(f') \cap Z'_{0} = \{x\}$. Le feuillet s'arrête sur cette
phrase.

L'exemple est donc complet dans sa conception et inachevé dans sa rédaction : il
donne tous les ingrédients et ne fait aucune des vérifications.

\subsection*{Une proposition sur trois algèbres graduées (pages 95 à 98)}

Quatre feuillets à l'encre bleue, mis au net et agrafés ensemble : ce sont les
seuls du dossier à avoir la forme d'un texte achevé, et ils traitent d'un sujet
qui n'a de rapport visible ni avec le paragraphe 12 ni avec le reste.

\emph{Proposition.} Soient $K$ un anneau, $A$, $B$, $C$ des $K$-algèbres
graduées à degrés positifs, réduites à $K$ en degré $0$, et
\[
  A \xrightarrow{\;\psi\;} B \xrightarrow{\;\varphi\;} C
\]
des homomorphismes. On suppose

\begin{itemize}
  \item[a)] pour tout $i \geqslant 0$, $\varphi^{i} : B^{i} \to C^{i}$ est
  surjectif et admet un inverse à droite $K$-linéaire $u^{i} : C^{i} \to B^{i}$.
\end{itemize}

\emph{(1)} Sous a), les deux conditions suivantes sont équivalentes :

\begin{itemize}
  \item[b)] $\mathrm{Ker}\,\varphi = \psi(A^{+})\,B$ ;
  \item[b')] pour tout $n \geqslant 0$, l'homomorphisme
  \[
    v^{n} : \bigoplus_{i+j=n} A^{i} \otimes_{K} C^{j} \longrightarrow B^{n} ,
    \qquad v^{n}(a^{i} \otimes c^{j}) = \psi(a^{i})\, u^{j}(c^{j}) ,
  \]
  est surjectif.
\end{itemize}

En particulier b') ne dépend pas du choix des $u^{j}$, ce que le feuillet note
entre crochets.

\emph{(2)} Si l'on suppose de plus

\begin{itemize}
  \item[b'')] $v^{n}$ est \emph{bijectif},
\end{itemize}

alors $\psi$ est injectif, et les $B^{i}$ sont plats — resp. projectifs — sur
$K$ si et seulement si les $A^{i}$ et les $C^{i}$ le sont.

\emph{(3)} Pour $x \in \mathrm{Spec}\,K$, posons
\[
  P_{A,x}(t) = \sum_{i} \mathrm{rg}_{k(x)}\bigl(A^{i}(x)\bigr)\, t^{i}
  \in \mathbf{Z}[[t]] ,
\]
et de même $P_{B,x}$, $P_{C,x}$. Alors on a, coefficient par coefficient,
\[
  P_{B,x} \ll P_{A,x}\, P_{C,x} ,
  \tag{$*$}
\]
avec égalité si b'') est vérifié ; et réciproquement, si les $A^{i}$, $B^{i}$,
$C^{i}$ sont plats de présentation finie — c'est-à-dire projectifs de type fini
— alors la condition

\begin{itemize}
  \item[c)] $P_{B,x} = P_{A,x}\, P_{C,x}$ pour tout $x$
\end{itemize}

implique b'').

C'est une proposition de \emph{scindage gradué} : sous a) et b''), $B$ est,
comme $K$-module gradué, le produit tensoriel $A \otimes_{K} C$, et la série de
Poincaré se multiplie. L'inégalité $(*)$ vaut toujours et l'égalité caractérise
le cas scindé — c'est la forme algébrique de l'énoncé qu'on connaît pour une
fibration, où la série de Poincaré de l'espace total est majorée par le produit
de celles de la base et de la fibre, avec égalité exactement lorsque la suite
spectrale dégénère.\footnote{Ce rapprochement est nôtre ; le dossier ne nomme
ni fibration ni suite spectrale, et n'écrit rien d'autre que le contenu
algébrique. La démonstration de $(*)$ n'est pas sur les feuillets : ils énoncent
l'inégalité et passent à l'égalité. La marge de la page 96, écrite en oblique
sur toute la hauteur, n'est lisible que par fragments ; le seul nom sûr qu'elle
porte est celui de Nakayama, invoqué pour le sens réciproque de (3), ce qui est
en effet l'argument : le rang d'un projectif de type fini se lit fibre à fibre,
et Nakayama relève la surjectivité.}

\emph{Corollaire 1 (page 97).} Moyennant a), les conditions suivantes sont
équivalentes : (i) $v^{*}$ bijectif et les $B^{i}$ projectifs de type fini ;
(ii) $v^{*}$ bijectif et les $A^{i}$, $C^{i}$ projectifs de type fini ; (iii)
les $A^{i}$, $B^{i}$, $C^{i}$ projectifs de type fini et $P_{B,x} = P_{A,x}
P_{C,x}$ pour tout $x$. Si $K$ est local, il suffit de regarder au point fermé.

Il nomme alors la situation : l'ensemble des conditions a), b'') sera dit
« bonnes conditions », et « très bonnes » si de plus les $B^{i}$ sont projectifs
de type fini.

\emph{Corollaire 2 (page 98).} Supposons $C \simeq K[t_{1},\ldots,t_{r}]$ avec
les $t_{i}$ homogènes de degrés $> 0$. Alors (a, b'') équivaut à
\[
  \begin{cases}
    B \simeq A[t_{1},\ldots,t_{r}] , \\
    \psi \simeq \text{l'homomorphisme canonique } A \to A[t_{1},\ldots,t_{r}] , \\
    \varphi \simeq \text{la projection } A[t_{1},\ldots,t_{r}] \to
      (A/A^{+})[t_{1},\ldots,t_{r}] .
  \end{cases}
\]
Le choix d'un tel isomorphisme équivaut à celui d'un homomorphisme d'algèbres
$u : C \to B$ avec $\varphi u = \mathrm{id}_{C}$, ou encore à celui d'un
homomorphisme $K$-linéaire relevant $B^{+} \to C^{+}/C^{+2}$. C'est le cas où le
scindage est réalisé par une section algébrique et non seulement linéaire, et il
dit exactement ceci : sous les bonnes conditions, une algèbre de polynômes en
quotient force $B$ à être une algèbre de polynômes sur $A$.

\emph{Question (page 97).} Les feuillets se ferment sur trois implications qu'il
pose et ne résout pas. Sous les « bonnes conditions », peut-on prouver :

\begin{enumerate}
  \item[(i)] $A$ plat sur $K$ $\Longrightarrow$ $C$ plat sur $B$ ;
  \item[(ii)] $C$ plat sur $K$ $\Longrightarrow$ $B$ plat sur $A$ ;
  \item[(iii)] $A$ et $C$ plats sur $K$ $\Longrightarrow$ $B$ plat sur $K$.
\end{enumerate}

La deuxième et la troisième sont celles que le scindage $B \simeq A \otimes_{K}
C$ rend plausibles. La première ne suit pas ce dessin.\footnote{« $C$ plat$/B$ »
est écrit ainsi sur le feuillet. $C$ étant un quotient de $B$, la platitude de
$C$ sur $B$ serait une condition très forte, et le parallélisme avec (ii)
ferait plutôt attendre « $C$ plat$/K$ ». Rien n'est corrigé ici : c'est une
\emph{question} et non une assertion, et une question mal posée reste une
question. Le feuillet, et la cote, s'arrêtent sur ces trois lignes.}

\end{document}
